\providecommand{\grpIrr}{\operatorname{Irr}}
\providecommand{\grpSt}{\mathrm{St}}
\providecommand{\grpFix}{\operatorname{Fix}}
\providecommand{\grpsgn}{\operatorname{sgn}}
\providecommand{\grpstd}{\mathrm{std}}
\providecommand{\grpone}{\mathbf{1}}
\providecommand{\grpDiff}{\mathfrak{D}}
\providecommand{\grprd}{\operatorname{rd}}
\providecommand{\grppos}[1]{#1^{+}}
\section{Nonabelian zero detection: general reductions}\label{sec:nab}
\providecommand{\nabFix}{\operatorname{Fix}}
\providecommand{\nabIrr}{\operatorname{Irr}}
\providecommand{\nabsgn}{\operatorname{sgn}}
\providecommand{\nabone}{\mathbf{1}}
\providecommand{\nabpos}[1]{(#1)^{+}}

The input theorem \cref{src:thm:dgrh} gives the Riemann hypothesis for the Dedekind zeta
function of every abelian extension of $\Q$. This section collects the general statements
that reduce the Riemann hypothesis for $\zeta_E$, with $E/\Q$ a nonabelian Galois extension
or an arbitrary number field, to a one-sided inequality for a sum over rational primes whose
weight is a class function of Frobenius. The reductions are equivalences. None of the
target inequalities is proved, and no new case of the Riemann hypothesis follows. We also determine, by finite group theory, which Frobenius
weights can and cannot serve in such a test, and how much of a nonabelian weight a fixed
congruence condition can see. The analytic interface of \cref{nab:sec:interface} is the
one used throughout Part~II. The worked symmetric-group instances (degree eight, and a
three-subgroup identity in degree five) are placed with the general results they
illustrate, in \cref{nab:sec:sn,nab:sec:reflection}. Detectors for particular groups and
fields are in \cref{sec:grp}, including the $A_5$ support frontier
(\cref{nab:prop:a5frontier}); the analytic engine is in \cref{sec:ana}, which also holds
the sampling bounds for positive ramp measures (\cref{nab:sec:sampling}).

\subsection{Setting and the analytic interface}\label{nab:sec:interface}
\begingroup\raggedright
\sources{grh_sylow_transfer_2026-11-25; grh_cyclic_generator_transfer_2026-11-25; grh_reflection_transfer_2026-11-25; grh_relative_prime_races_2026-11-25; grh_deficit_transfer_2026-11-25; grh_method_transfers_2026-11-25; grh_a6_transfer_2026-11-25; psl2_grh_transfer; grh-transfer}
\endgroup

Throughout, $E/\Q$ is a finite Galois extension with group $G$, all fields are fixed, and no
uniformity in the degree or the discriminant is claimed. We write $S$ for a finite set of
rational primes containing every prime ramified in $E$, and $g_p\in G$ for a Frobenius
element at $p\notin S$; the value $w(g_p)$ of a class function $w$ does not depend on the
choice. For a subgroup $U\le G$ the permutation character of $G/U$ is
\[
 \Ind_U^G\nabone(g)=\#\nabFix(g,G/U)=\frac{\abs{\set{x\in G: x^{-1}gx\in U}}}{\abs{U}},
\]
and at $p\notin S$ its value at $g_p$ is the number $a_K(p)$ of primes of degree one above
$p$ in $K=E^U$. If $K=\Q[T]/(f)$ with $f\in\Z[T]$ monic irreducible and $p\nmid\disc(f)$,
then $a_K(p)$ is the number of roots of $f$ in $\F_p$. When a weight is given by the
factorization pattern of a polynomial modulo $p$, the primes dividing its discriminant are
placed in $S$.

\paragraph{Prime sums.}
For a real class function $w$ on $G$ with mean $\delta_w=\langle w,\nabone\rangle$ put
\begin{align*}
 P_w(x)&=\sum_{\substack{p\le x\\ p\notin S}}w(g_p)\log p\Bigl(\frac1p-\frac1x\Bigr)\quad(x\ge1),
 &
 \theta_w(x)&=\sum_{\substack{p\le x\\ p\notin S}}w(g_p)\log p,\\
 D^{\mathrm{pr}}_w(s)&=\sum_{p\notin S}\frac{w(g_p)\log p}{p^{s}}\quad(\Re s>1). &&
\end{align*}
In the notation of \cref{ana:sec:primeonly}, $P_w$ is the prime-only ramp, $\theta_w$ is
$\Theta_w$ and $\delta_w$ is $a_w$; the Dirichlet series $D^{\mathrm{pr}}_w$ is the
prime-only part of the series $D_w$ there. The analytic engine is \cref{ana:thm:primeonly},
the prime-only form of the one-sided Mellin--Landau test. Its hypothesis is a
representation of the prime-power series of $w$ by logarithmic derivatives of products of
Hecke $L$-functions with positive weights; the verification of that hypothesis is the
group-theoretic part. We recall, in the present notation, the part of it used in Part~II
(\cref{sec:nab,sec:grp,sec:hec}), and then its two instances used below: an entire product
of induced characters (\cref{grp:cor:detector}) and Dedekind zeta functions
(\cref{nab:prop:dedekind}).

\begin{proposition}[Analytic interface; \stU]\label[proposition]{grp:prop:interface}
With the notation above, the following hold.
\begin{enumerate}[label=\textup{(\alph*)}]
\item \textup{(\eqref{ana:eq:Rw} and \cref{ana:lem:mellin}.)} The limit
$c_w=\lim_{x\to\infty}\bigl(P_w(x)-\delta_w\log x\bigr)$ exists. The centred
remainder
\[
 R_w(x)=P_w(x)-\delta_w\log x-c_w
\]
is continuous and bounded on $[1,\infty)$ and tends to $0$, and for $\Re s>1$
\begin{equation}\label{grp:eq:mellin}
 \int_1^\infty R_w(x)x^{-s}\,\dd x
 =\frac{D^{\mathrm{pr}}_w(s)}{s(s-1)}-\frac{\delta_w}{(s-1)^2}-\frac{c_w}{s-1}.
\end{equation}
\item \textup{(\cref{ana:thm:primeonly,ana:thm:zerotest}.)} Suppose that for
$\Re s>1$
\begin{equation}\label{grp:eq:decomp}
 D^{\mathrm{pr}}_w(s)=-\sum_{j=1}^{r}a_j\frac{Q_j'(s)}{Q_j(s)}+H(s),
\end{equation}
where the $a_j$ are positive, every $Q_j$ is a finite product of integer powers
of $L$-functions of finite-order Hecke characters of number fields, Dedekind
zeta functions included, with no pole in $1/2<\Re s<1$, and $H$ is holomorphic
on $\Re s>1/2$ and bounded on $\Re s\ge1/2+\eta$ for every $\eta>0$. Then the
poles of \eqref{grp:eq:mellin} in $1/2<\Re s<1$ are simple and lie at the zeros
$\rho$ of the $Q_j$, with residue
\begin{equation}\label{grp:eq:residue}
 \frac{\sum_j a_j\ord_\rho Q_j}{\rho(1-\rho)} ,
\end{equation}
which is positive when $\rho$ is real. If $1/2<\beta<1$ and
$R_w(x)\le Cx^{\beta-1}$ for all sufficiently large $x$, then no $Q_j$ has a
zero in $\Re s>\beta$; in particular an eventual bound $R_w(x)\le Cx^{-1/2}$
excludes the zeros of every $Q_j$ in $\Re s>1/2$.
\item \textup{(\cref{ana:thm:primeonly}\textup{(iv)} and \cref{ana:rem:finite}.)}
In the situation of \textup{(b)}, suppose that for every $0<\eps<1/2$ there is
$K_\eps$ with
\begin{equation}\label{grp:eq:dyadic}
 P_w(n)-P_w(2n)+\delta_w\log2\le K_\eps n^{-1/2+\eps}
\end{equation}
for all sufficiently large integers $n$. Then
$R_w(x)\le K'_\eps x^{-1/2+\eps}$ for all large real $x$, and no $Q_j$ has a
zero in $\Re s>1/2$. The left side of \eqref{grp:eq:dyadic} equals
\[
 R_w(n)-R_w(2n)=-\frac{\theta_w(n)}{2n}
 -\sum_{n<p\le2n,\ p\notin S}w(g_p)\log p\Bigl(\frac1p-\frac1{2n}\Bigr)+\delta_w\log2,
\]
so no centring constant enters.
\item \textup{(\cref{ana:thm:primeonly}\textup{(iii)} and its last assertion.)}
In the situation of \textup{(b)}, if no $Q_j$ has a zero in $\Re s>1/2$, then
$R_w(x)=O_\eps(x^{-1/2+\eps})$ for every $\eps>0$; if $\zeta_E$ satisfies the
Riemann hypothesis, then $R_w(x)=O_{E,w}(x^{-1/2}\log^2(2x))$.
\end{enumerate}
\end{proposition}

The prime-only series $D^{\mathrm{pr}}_w$ differs from $D_w$ by a series that converges
normally on $\Re s\ge1/2+\eta$ for every $\eta>0$ (\cref{ana:lem:powers}), so
\eqref{grp:eq:decomp} is the hypothesis of \cref{ana:thm:primeonly}, and the
residue \eqref{grp:eq:residue} is \eqref{ana:eq:residue}, through
\cref{ana:lem:mellin}\textup{(iv)}. The kernel identity behind
\eqref{grp:eq:mellin} is \eqref{ana:eq:mellin} for the prime-only coefficients.

\begin{remark}[Variants]\label[remark]{grp:rem:variants}
Three variants of \cref{grp:prop:interface} are used below, all owned by
\cref{sec:ana}.
First, one may sum over all prime powers and all primes, with the coefficient
$a_J(p^k)=\log p\,\tr\bigl(g_p^k\mid V_J^{I_p}\bigr)
=\log p\,\abs{I_p}^{-1}\sum_{h\in I_p}J(g_p^kh)$ of
$-L'(s,J)/L(s,J)$ for a representation $V_J$ with character $J$, inertia group
$I_p$ and Frobenius lift $g_p$, and with the constant
$c_J=b_J-\delta_J$ read from
$-L'/L(s,J)=\delta_J/(s-1)+b_J+O(s-1)$ (\cref{ana:def:ramp}). By
\cref{ana:lem:powers} and \eqref{ana:eq:RwR} the two remainders differ by
$O(x^{-1/2})$, the constants being related by $c_{J,S}=c_J-\sum a_J(n)/n$, the
sum over all omitted prime powers; the criterion at the scale
$x^{-1/2+\eps}$ is unchanged.
Second, for the full prime-power remainder $R$ the pointwise bound may be
replaced by the mean bound $\int_X^{2X}\grppos{R(x)}\,\dd x=O_\eps(X^{\beta+\eps})$
for the positive part (\cref{ana:thm:zerotest,ana:cor:forms}).
Third, the difference $P_w(n)-P_w(n^2)+\delta_w\log n$ at the integer scales
$n,n^2$ may replace the dyadic difference in \textup{(c)}
(\cref{ana:thm:primeonly}\textup{(v)}).
\end{remark}

\begin{corollary}[Monomially certified detectors; \stU]\label[corollary]{grp:cor:detector}\label{nab:prop:signed}
Let $J=a\,\nabone+\sum_i\Ind_{H_i}^G\lambda_i$ with an integer $a\ge0$ and
nontrivial linear characters $\lambda_i$ of subgroups $H_i\le G$, and assume
that $J$ is real-valued. Then $L(s,J)$ is holomorphic on $\C\setminus\set1$,
has a pole of order $a=\delta_J$ at $1$, and does not vanish on $\Re s\ge1$.
\begin{enumerate}[label=\textup{(\roman*)}]
\item If $1/2\le\beta<1$ and $R_J(x)\le Cx^{\beta-1}$ for all large $x$, then
$L(s,J)\ne0$ for $\Re s>\beta$.
\item Suppose in addition that there is a constant $k$ with
$\ord_\rho\zeta_E\le k\,\ord_\rho L(s,J)$ for all $\rho$ with $1/2<\Re\rho<1$.
Then the following are equivalent: the Riemann hypothesis for $\zeta_E$; for
every $\eps>0$, $R_J(x)\le K_\eps x^{-1/2+\eps}$ for all large $x$; for every
$\eps>0$, $R_J(x)=O_\eps(x^{-1/2+\eps})$; the dyadic condition
\eqref{grp:eq:dyadic} for $w=J$. If they hold, every Artin $L$-function of
$E/\Q$ is holomorphic and nonzero in $1/2<\Re s<1$.
\end{enumerate}
\end{corollary}

\begin{proof}
The first assertion is Hecke's theorem for the factors $L(s,\lambda_i,E/E^{H_i})$.
For unramified $p$ the coefficient of $p^{-ks}\log p$ in $-L'/L(s,J)$ is
$J(g_p^k)$, of modulus at most $J(1)$; the terms with $k\ge2$ converge
normally on $\Re s\ge1/2+\eta$ for every $\eta>0$ and the finitely many local
series at $p\in S$ converge normally on $\Re s>0$ (\cref{ana:lem:powers}).
Hence $D^{\mathrm{pr}}_J=-L'/L(s,J)-H_J$ with $H_J$ holomorphic on $\Re s>1/2$ and
bounded on each half-plane $\Re s\ge1/2+\eta$, which is \eqref{grp:eq:decomp} with
$Q_1=L(s,J)$, a finite product of Hecke $L$-functions. Now (i) is
\cref{grp:prop:interface}(b). In (ii), zero-freeness of $L(s,J)$ on
$\Re s>1/2$ gives that of $\zeta_E$; the functional equation and the classical
nonvanishing on $\Re s=1$ give the Riemann hypothesis. The converse and the
two-sided bound are \cref{grp:prop:interface}(d), the dyadic form is
\cref{grp:prop:interface}(c), and the last assertion is
Lemma~\ref{nab:lem:heilbronn}(c).
\end{proof}

The case $a=0$ is the one used in \cref{nab:sec:sylow}. Then
$J=\sum_i\Ind_{H_i}^G\lambda_i$ with every $\lambda_i$ nontrivial, $L(s,J)$ is entire,
$\delta_J=\langle J,\nabone\rangle=0$, and $R_J=P_J-c_J$ with $c_J=\lim_{x\to\infty}P_J(x)$.
Part (i) and \cref{grp:prop:interface}(d), with $Q_1=L(s,J)$, give: $L(s,J)\ne0$ for
$\Re s>1/2$ if and only if, for every $\eps>0$, $R_J(x)\le C_\eps x^{-1/2+\eps}$ for all
sufficiently large $x$, and then $R_J(x)=O_\eps(x^{-1/2+\eps})$; under the Riemann
hypothesis for $\zeta_E$, $R_J(x)=O_{E,J}(x^{-1/2}\log^2(2x))$.

Replacing an inequality of \cref{grp:cor:detector} by its negative does not
give the same conclusion, because the sign of the residue
\eqref{grp:eq:residue} at a real pole is reversed.

\begin{proposition}[Positive weights and Dedekind zeta functions]\label{nab:prop:dedekind}
Let $K\subseteq E$ be a subfield. There is a unique constant $c_K=c_{K,S}$ for which
\[
 R_K(x)=\sum_{\substack{p\le x\\p\notin S}}a_K(p)\log p\Bigl(\frac1p-\frac1x\Bigr)-\log x-c_K
\]
tends to $0$ as $x\to\infty$. Equivalently, with
$\theta_K(u)=\sum_{p\le u,\,p\notin S}a_K(p)\log p$, one has
$R_K(x)=-\int_x^\infty(\theta_K(u)-u)u^{-2}\,\dd u$.
\begin{enumerate}[label=\textup{(\alph*)}]
\item Let $K_1,\dots,K_q\subseteq E$ and $w_1,\dots,w_q>0$. If
$\sum_jw_jR_{K_j}(x)\le Cx^{\beta-1}$ for some $1/2<\beta<1$ and all sufficiently large
$x$, then no $\zeta_{K_j}$ has a zero in $\Re s>\beta$.
\item If the Riemann hypothesis holds for $\zeta_K$, then $R_K(x)=O_K(x^{-1/2})$.
\end{enumerate}
\end{proposition}
\begin{proof}
The existence of $c_K$ and part (a) are \cref{ana:thm:primeonly} for the class function
$w=\sum_jw_j\Ind_{U_j}^G\nabone$, $K_j=E^{U_j}$, with the functions $F_j=\zeta_{K_j}$ and
$a_w=\sum_jw_j$ (\cref{ana:ex:instances}(a)); its remainder is $R_w=\sum_jw_jR_{K_j}$, since
both tend to $0$. For (b), under the Riemann hypothesis
the integrated explicit formula gives $\int_0^x(\psi_K(t)-t)\,\dd t=O_K(x^{3/2})$
\cite{GrenieMolteni2016Smoothed}, and $\psi_K(t)-\theta_K(t)=O_K(\sqrt t)$ by the Chebyshev
bound for squares, the bound $O_K(t^{1/3}\log(2t))$ for higher powers and $O_K(\log(2t))$
for the primes of $S$. Hence $A(x)=\int_1^x(\theta_K(t)-t)\,\dd t=O_K(x^{3/2})$, and
integration by parts gives $R_K(x)=A(x)x^{-2}-2\int_x^\infty A(t)t^{-3}\,\dd t=O_K(x^{-1/2})$.
\end{proof}

\begin{remark}\label{grp:prop:split}
In (a) with $q=1$ and the exact exponent, an eventual bound $R_K(x)\le Cx^{-1/2}$ implies
the bound for every $\beta>1/2$, hence the Riemann hypothesis for $\zeta_K$ by the
functional equation. Thus $\mathrm{RH}(\zeta_K)$ is equivalent to $R_K(x)\le Cx^{-1/2}$
for large $x$, and also to $R_K(x)\le C_\eps x^{-1/2+\eps}$ for every $\eps>0$. For
Galois $K=M$ of degree $n$ the centred ramp over completely split primes used in
\cref{grp:sec:kummer} is $R_M/n$; the exact value of $\limsup\sqrt x\,R_M(x)/n$ under
$\mathrm{RH}(\zeta_M)$ is \cref{zer:thm:endpoint}.
\end{remark}

Two further forms of the test are used. For a class function $w$ of mean zero whose
prime-power series is a positive combination of logarithmic derivatives of entire quotients
of Dedekind zeta functions, \cref{ana:thm:primeonly}(i)$\Leftrightarrow$(v) is the
constant-free integer two-scale criterion: no quotient has a zero in $\Re s>1/2$ if and only
if, for every $0<\eps<1/2$, $P_w(x)-P_w(x^2)\le C_\eps x^{-1/2+\eps}$ for all sufficiently
large integers $x$; no logarithmic term appears because $\langle w,\nabone\rangle=0$. The
second is the averaged finite-difference detector \cref{ana:thm:fd}. For a fixed $c>1$ its
statistic $T(cx)-T(x)$ has the nonnegative kernel
\[
 K_c(y)=1-c^{-1}\ (0<y\le1),\qquad K_c(y)=y^{-1}-c^{-1}\ (1<y\le c),\qquad K_c(y)=0\ (y>c),
\]
and its Mellin multiplier $H_c(s)=(c^{s-1}-1)/(s(s-1))$ is positive on the positive real
axis, whereas the kernel $1/(s(s-1))$ of the ramp $1/n-1/x$ is negative on $(0,1)$
(\cref{ana:lem:fd}). This is why a nonnegative weight is tested by an upper bound for a
finite difference in \cref{nab:thm:deficit} and by a lower bound for the ramp in
\cref{nab:thm:nonsplit}.

\subsection{Heilbronn characters and multiplicity domination}\label{nab:sec:heilbronn}
\begingroup\raggedright
\sources{grh_sylow_transfer_2026-11-25; grh_cyclic_generator_transfer_2026-11-25; grh_reflection_transfer_2026-11-25; grh_method_transfers_2026-11-26; psl2_grh_transfer; grh_deficit_transfer_2026-11-25}
\endgroup

Let $E/k$ be a Galois extension of number fields with group $G$ and let $s_0\in\C$,
$s_0\ne1$. By Brauer induction every Artin $L$-function $L(s,\chi,E/k)$ is meromorphic.
Its order $n_\chi=\ord_{s=s_0}L(s,\chi,E/k)$ is therefore an integer, negative at a pole. The
\emph{Heilbronn character} at $s_0$ is the virtual character
\[
 \Theta=\Theta_{s_0}=\sum_{\chi\in\nabIrr(G)}n_\chi\chi,\qquad d=\Theta(1)=\ord_{s_0}\zeta_E .
\]

\begin{lemma}[Heilbronn--Stark facts; \stU]\label{nab:lem:heilbronn}
Let $\Theta$ be the Heilbronn character of $E/k$ at a point $s_0\ne1$.
\begin{enumerate}[label=\textup{(\alph*)}]
\item For every subgroup $U\le G$ and every character $\eta$ of $U$,
$\langle\Res_U\Theta,\eta\rangle=\ord_{s_0}L(s,\eta,E/E^U)$. In particular
$\langle\Theta,\Ind_U^G\nabone\rangle=\ord_{s_0}\zeta_{E^U}$.
\item If every irreducible character of $U$ is induced from a linear character of a
subgroup, then $\Res_U\Theta$ is the character of an actual representation of $U$ of
dimension $d$, possibly zero. This holds for cyclic, abelian and dihedral subgroups and for
subgroups of prime-power order. In particular $d\ge0$.
\item If $d=0$, then $\Theta=0$: every irreducible Artin $L$-function of $E/k$ is
holomorphic and nonzero at $s_0$.
\end{enumerate}
\end{lemma}
\begin{proof}
(a) is Frobenius reciprocity together with the invariance of Artin $L$-functions under
induction and their multiplicativity on virtual characters; complex irreducible characters
therefore create no sign ambiguity in the integer orders. (b) A monomial irreducible
character $\eta=\Ind_V^U\lambda$ has $L(s,\eta,E/E^U)=L(s,\lambda;E^V)$, a Hecke
$L$-function of a character of finite order, which is holomorphic at $s_0\ne1$; this
includes the trivial character, whose Dedekind zeta function has its only pole at $1$.
Hence all multiplicities of $\Res_U\Theta$ are nonnegative. Groups of prime-power order
are monomial; for a dihedral group the irreducible characters are linear or induced from
the cyclic rotation subgroup. (c) For every $g\in G$ the restriction of $\Theta$ to
$\langle g\rangle$ is the character of a representation of dimension $0$, so
$\Theta(g)=0$. These are the classical facts of Heilbronn and Stark
\cite{Heilbronn1973Real,FooteGinsbergMurty2015Heilbronn}. No holomorphy of a nonmonomial
Artin $L$-function is used, and $\Theta$ itself is not asserted to be an actual character.
\end{proof}

Part (c) is used repeatedly: once the Riemann hypothesis is known for $\zeta_E$, every
irreducible Artin $L$-function of $E/\Q$ is holomorphic and nonzero in $1/2<\Re s<1$. This
is zero and pole exclusion in the open half-strip. It says nothing about poles at zeros on
the critical line, so it does not give Artin holomorphy.

Now let $k=\Q$ and fix $\rho$ with $0<\Re\rho<1$. Let $H\le G$ be a nontrivial subgroup,
$B=E^H$ and $b=\ord_\rho\zeta_B$, so that $b=\langle\Theta,\Ind_H^G\nabone\rangle$ by
Lemma~\ref{nab:lem:heilbronn}(a). For each prime $\ell\mid\abs H$ choose a Sylow
$\ell$-subgroup $P_\ell\le H$, and let $r$ be the number of primes dividing $\abs H$. For
a cyclic subgroup $C\le G$ put
\[
 W_C=\Reg_G-\Ind_C^G\nabone=\Ind_C^G(\Reg_C-\nabone),\qquad
 D(g)=\langle\Theta,W_{\langle g\rangle}\rangle=\ord_\rho\frac{\zeta_E}{\zeta_{E^{\langle g\rangle}}} .
\]
The cyclic-generator weight is defined for $H=[G,G]$ only. When $G$ is nonabelian and
$H=[G,G]$, choose for each $\ell\mid\abs H$ a finite list $S_\ell$ of nonidentity elements
generating $P_\ell$, let $\mathcal T$ contain one representative of each $G$-conjugacy
class of cyclic subgroups $\langle g\rangle$, $g\in\bigcup_\ell S_\ell$, and let $M$ be the
largest number of elements of the combined lists $S_\ell$ generating conjugate cyclic
subgroups. Define the characters
\begin{equation}\label{nab:eq:weights-def}
 W^{\mathrm{Syl}}=\sum_{\ell\mid\abs H}\bigl(\Reg_G-\Ind_{P_\ell}^G\nabone\bigr),
 \qquad
 W^{\mathrm{cyc}}=\sum_{C\in\mathcal T}W_C\quad(H=[G,G]).
\end{equation}
Each summand is induced from the augmentation character $\Reg_P-\nabone$ of a group $P$ of
prime-power order or of a cyclic group. Since $P$ is monomial and an irreducible
constituent of $\Reg_P-\nabone$ cannot be induced from a trivial character, both
$W^{\mathrm{Syl}}$ and $W^{\mathrm{cyc}}$ are sums of characters induced from nontrivial
linear characters, as required in the case $a=0$ of \cref{grp:cor:detector}. Their $L$-functions are
\[
 L(s,W^{\mathrm{Syl}})=\prod_{\ell\mid\abs H}\frac{\zeta_E(s)}{\zeta_{E^{P_\ell}}(s)},
 \qquad
 L(s,W^{\mathrm{cyc}})=\prod_{C\in\mathcal T}\frac{\zeta_E(s)}{\zeta_{E^{C}}(s)} .
\]

\begin{theorem}[Multiplicity domination; \stU]\label{nab:thm:heilbronn}
Let $\Theta$ be the Heilbronn character of $E/\Q$ at $\rho$, so that
$\langle\Res_U\Theta,\eta\rangle=\ord_\rho L(s,\eta,E/E^U)$ for every subgroup $U\le G$ and
every character $\eta$ of $U$, $\Res_U\Theta$ is the character of an actual
representation whenever $U$ is cyclic or of prime-power order, and $\Theta=0$ if
$d=\ord_\rho\zeta_E=0$ \textup{(Lemma~\ref{nab:lem:heilbronn})}. With the notation above:
\begin{enumerate}[label=\textup{(\alph*)}]
\item \textup{(Sylow form, any nontrivial $H\le G$.)} Let
$e_\ell=\ord_\rho(\zeta_E/\zeta_{E^{P_\ell}})$ and
$\delta_\ell=\#\set{h\in H:\ell\mid\ord(h)}/\abs H$. Then $0\le e_\ell\le d$ and
\begin{equation}\label{nab:eq:dom}
 d-b\;\le\;2\sum_{\ell\mid\abs H}\delta_\ell e_\ell\;\le\;2\,\ord_\rho L(s,W^{\mathrm{Syl}})
 \;\le\;2r\,d .
\end{equation}
If $b=0$ and $\ord_\rho L(s,W^{\mathrm{Syl}})=0$, then $d=0$ and every irreducible Artin
$L$-function of $E/\Q$ is holomorphic and nonzero at $\rho$.
\item \textup{(Cyclic-generator form, $G$ nonabelian and $H=[G,G]$.)} If $b=0$, then
$e_\ell\le\sum_{g\in S_\ell}D(g)$ for every $\ell$, and
\begin{equation}\label{nab:eq:domcyc}
 d\;\le\;2\sum_{\ell\mid\abs H}\sum_{g\in S_\ell}D(g)\;\le\;2M\,\ord_\rho L(s,W^{\mathrm{cyc}}) .
\end{equation}
If moreover $\ord_\rho L(s,W^{\mathrm{cyc}})=0$, then every irreducible Artin $L$-function
of $E/\Q$ is holomorphic and nonzero at $\rho$.
\end{enumerate}
\end{theorem}
\begin{proof}
(a) By Lemma~\ref{nab:lem:heilbronn}, $\Res_{P_\ell}\Theta$ is the character of a
representation $V_\ell$ of $P_\ell$ of dimension $d$, which may be taken unitary, and the
dimension of $V_\ell^{P_\ell}$ is $\ord_\rho\zeta_{E^{P_\ell}}$. Hence $e_\ell$ is the
codimension of the fixed subspace, so $0\le e_\ell\le d$.

Let $h\in H$ and write $h=\prod_{\ell\mid\ord(h)}h_\ell$ for its commuting components of
prime-power order in $\langle h\rangle$. Each $h_\ell$ is conjugate in $H$ to an element
$h_\ell'$ of $P_\ell$. In the actual representation of $\langle h\rangle$ with character
$\Res_{\langle h\rangle}\Theta$, the fixed space of $h_\ell$ has dimension
$\abs{\langle h_\ell\rangle}^{-1}\sum_{v\in\langle h_\ell\rangle}\Theta(v)$, which is
invariant under conjugation and so equals the dimension of the fixed space of $h_\ell'$
in $V_\ell$. The latter contains $V_\ell^{P_\ell}$, so the fixed space of $h_\ell$ has
codimension at most $e_\ell$. The fixed space of $h$ is the intersection of those of the
$h_\ell$, because each component is a power of $h$; its codimension is at most
$\sum_{\ell\mid\ord(h)}e_\ell$. The eigenvalues of $h$ are roots of unity, so each
eigenvalue different from $1$ contributes at most $2$ to $d-\Re\Theta(h)$:
\begin{equation}\label{nab:eq:pointdom}
 0\le d-\Re\Theta(h)\le2\sum_{\ell\mid\ord(h)}e_\ell .
\end{equation}
Average over $h\in H$. By Lemma~\ref{nab:lem:heilbronn}(a) the average of $\Theta(h)$ is
$\langle\Res_H\Theta,\nabone\rangle=b$, whether or not $\Res_H\Theta$ is an actual
character. This gives the first inequality of \eqref{nab:eq:dom}. The second holds because
$\delta_\ell\le1$ and $\sum_\ell e_\ell=\ord_\rho L(s,W^{\mathrm{Syl}})$, and the third
because $e_\ell\le d$. If $b=0$ and $\ord_\rho L(s,W^{\mathrm{Syl}})=0$, then $d\le0$, and
Lemma~\ref{nab:lem:heilbronn}(c) applies.

(b) Now $H=[G,G]$ and $b=0$. For $g\in S_\ell$ the cyclic group $\langle g\rangle$ lies in
$P_\ell$, so $D(g)=d-\dim V_\ell^{\langle g\rangle}\ge0$. Since $V_\ell^{P_\ell}$ is the
intersection of the fixed spaces of the generators in $S_\ell$, subadditivity of
codimension gives $e_\ell\le\sum_{g\in S_\ell}D(g)$. Averaging \eqref{nab:eq:pointdom}
over $h\in H$ with $b=0$ gives $d\le2\sum_\ell e_\ell$, hence the first inequality of
\eqref{nab:eq:domcyc}. Conjugate cyclic subgroups have equal $D(g)$, so the double sum is at
most $M\sum_{C\in\mathcal T}\langle\Theta,W_C\rangle=M\ord_\rho L(s,W^{\mathrm{cyc}})$. If
this order is zero, then $d=0$, and Lemma~\ref{nab:lem:heilbronn}(c) applies.
\end{proof}

Part (a) is the Sylow form, for an arbitrary nontrivial subgroup $H$; part (b) is the
cyclic-generator form, for $H=[G,G]$ under the hypothesis $b=0$. The second refines the
support of the first and is not a formal consequence of its statement, so both are
recorded.

\subsection{The reduction for an arbitrary Galois group}\label{nab:sec:sylow}
\begingroup\raggedright
\sources{grh_sylow_transfer_2026-11-25; grh_cyclic_generator_transfer_2026-11-25; psl2_grh_transfer}
\endgroup

\begin{theorem}[Sylow-filtered and cyclic-generator prime race; \stG; \stR]\label{nab:thm:sylow}
Let $E/\Q$ be a finite Galois extension with nonabelian group $G$, let $H=[G,G]$, and let
$B=E^H$ be the maximal abelian subfield of $E$; by \cref{src:thm:dgrh}, $\zeta_B$ satisfies
the Riemann hypothesis. Let $W^{\mathrm{Syl}}$ and $W^{\mathrm{cyc}}$ be the characters of
\eqref{nab:eq:weights-def} for this $H$, and for $W$ either of them let $R_W=P_W-c_W$ be as
in \cref{grp:prop:interface}(a), with $\delta_W=0$. For either choice of $W$ the following are
equivalent:
\begin{enumerate}[label=\textup{(\roman*)}]
\item the Riemann hypothesis for $\zeta_E$;
\item $L(s,W)$ has no zeros in $\Re s>1/2$;
\item for every $\eps>0$ there is $C_\eps$ with $R_W(x)\le C_\eps x^{-1/2+\eps}$ for all
sufficiently large $x$.
\end{enumerate}
If these hold, every irreducible Artin $L$-function of $E/\Q$ is holomorphic and nonzero in
$\Re s>1/2$, apart from the pole of the trivial one at $s=1$. For $W=W^{\mathrm{cyc}}$, if
the bound in \textup{(iii)} holds for a single $0<\eps<1/2$, then $\zeta_E$ has no zero,
and no irreducible Artin $L$-function of $E/\Q$ has a zero or a pole, in
$1/2+\eps<\Re s<1$. At every $\rho$ with $1/2<\Re\rho<1$,
\begin{equation}\label{nab:eq:mult}
 \tfrac12\ord_\rho\zeta_E\le\ord_\rho L(s,W^{\mathrm{Syl}})\le r\,\ord_\rho\zeta_E,
 \qquad
 \ord_\rho\zeta_E\le2M\,\ord_\rho L(s,W^{\mathrm{cyc}}),
\end{equation}
so $L(s,W)$ and $\zeta_E$ have the same zeros there.
\end{theorem}
\begin{proof}
At $\rho$ with $1/2<\Re\rho<1$ the Riemann hypothesis for $\zeta_B$ gives $b=0$, and
\eqref{nab:eq:mult} is \cref{nab:thm:heilbronn}(a),(b). Both $L$-functions are products of
quotients $\zeta_E/\zeta_{E^P}$, with $P$ a Sylow subgroup $P_\ell$ or a cyclic subgroup
conjugate to a subgroup of some $P_\ell$, and by the proof of \cref{nab:thm:heilbronn} the order of such a quotient at
$\rho$ is the codimension of the $P$-fixed space in a representation of dimension $d$,
hence lies between $0$ and $d$. So a zero of $L(s,W)$ at $\rho$ is a zero of $\zeta_E$, and
$L(s,W)$ and $\zeta_E$ have the same zeros in $1/2<\Re s<1$. Both are nonzero on $\Re s\ge1$
apart from the pole of $\zeta_E$, and the functional equation of $\zeta_E$ gives the left
half of the strip. This proves (i)$\Leftrightarrow$(ii). The equivalence with (iii) is
\cref{grp:cor:detector}(i) and \cref{grp:prop:interface}(d), with $Q_1=L(s,W)$. For
$W^{\mathrm{cyc}}$ and a single $\eps$, \cref{grp:cor:detector}(i) with $\beta=1/2+\eps$
excludes the zeros of
$L(s,W^{\mathrm{cyc}})$ in $\Re s>1/2+\eps$, and \eqref{nab:eq:mult} those of $\zeta_E$. The
assertions on Artin $L$-functions in the strip are Lemma~\ref{nab:lem:heilbronn}(c); on
$\Re s\ge1$ they are classical.
\end{proof}

Under (i), the last assertion of \cref{grp:prop:interface}(d), which rests on the
effective Chebotarev theorem under the Riemann hypothesis for $\zeta_E$, gives the sharper
bound $R_W(x)=O(x^{-1/2}\log^2(2x))$. For $W^{\mathrm{Syl}}$ the restriction to $H=[G,G]$ only
serves to make the hypothesis on $\zeta_B$ known: for every nontrivial subgroup $H\le G$,
with $W^{\mathrm{Syl}}$ built from Sylow subgroups of $H$ and the Riemann hypothesis for
$\zeta_{E^H}$ assumed, (i)--(iii) are equivalent, by the same proof, which uses only
\cref{nab:thm:heilbronn}(a) for $W^{\mathrm{Syl}}$.

\paragraph{Which primes enter.}
For a subgroup $P\le G$ and a cyclic subgroup $C\le G$,
\begin{equation}\label{nab:eq:weights}
 \bigl(\Reg_G-\Ind_P^G\nabone\bigr)(g)=
 \begin{cases}\abs G-[G:P],&g=1,\\ -\#\nabFix(g,G/P),&g\ne1,\end{cases}
 \qquad
 W_C(u)=-\frac{\abs{C_G(u)}\,\abs{C\cap u^G}}{\abs C}\quad(u\ne1).
\end{equation}
Both weights are positive only at the identity and have mean zero over $G$, so $R_W$ has
no logarithmic main term. For $W^{\mathrm{Syl}}$ a nonidentity $g$ has nonzero weight only
if it is conjugate to an element of prime-power order of $H$; since $H=[G,G]$ is normal,
this says that $g\in H$ has prime-power order. For $W^{\mathrm{cyc}}$ the negative support
consists of the classes meeting the chosen cyclic subgroups; when every $P_\ell$ is
generated by elements of order $\ell$, only Frobenius elements of prime order enter. The
test ignores every prime that does not split completely in the maximal abelian subfield
$B$: this is exactly where Dirichlet GRH is used. One may add a constant to make the
weights nonnegative; the centered sum of the constant is of the required size by the
Riemann hypothesis for $\zeta_\Q$, so the criterion is unchanged.

\begin{remark}\label{nab:rem:scope}
The test is a centered, signed race attached to an entire product of induced Hecke
$L$-functions. It is not claimed to have least support: the classical criterion obtained
from $-\zeta_E'/\zeta_E$, with the pole removed and prime powers controlled, tests
$\mathrm{RH}(\zeta_E)$ using completely split primes only
(Proposition~\ref{nab:prop:dedekind} with $K=E$). The signed family removes the abelian quotient, has
zero main term, and supplies the multiplicity bound \eqref{nab:eq:mult}. Nothing here
shows that its arithmetic estimate is easier.
\end{remark}

\begin{corollary}[Arbitrary number fields and Hecke characters of finite order; \stG; \stR]\label{nab:cor:any}
Let $K$ be a number field whose Galois closure $E$ over $\Q$ has nonabelian group $G$. The
single estimate \textup{(iii)} of \cref{nab:thm:sylow} for $E$, which concerns $H=[G,G]$,
implies the Riemann hypothesis for $\zeta_K$. If $L/K$ is the abelian extension attached by class field theory
to a Hecke character $\chi$ of finite order of $K$, and \textup{(iii)} holds for the Galois
closure of $L$ over $\Q$, then $L(s,\chi)$ has no zeros in $\Re s>1/2$.
\end{corollary}
\begin{proof}
By \cref{nab:thm:sylow} the estimate gives the Riemann hypothesis for $\zeta_E$. The
extension $E/K$ is Galois, so $\zeta_E/\zeta_K$ is entire, by the Aramata--Brauer
theorem \cite{Brauer1947Zeta} or by Lemma~\ref{nab:lem:heilbronn}(a) applied to the
trivial character of $\Gal(E/K)$. A zero of $\zeta_K$ in $1/2<\Re s<1$ is therefore a zero of $\zeta_E$, and the
functional equation completes the argument. For the second statement, $\zeta_L$ satisfies
the Riemann hypothesis by the first, and $\zeta_L=\prod_\chi L(s,\chi)$ over the characters
of $\Gal(L/K)$ is a product of functions holomorphic in the right half of the strip.
\end{proof}

The second statement concerns characters of finite order only. It gives no information on
Hecke characters of infinite order or on general automorphic $L$-functions.

\begin{remark}[Abelian base fields]\label{nab:rem:basechange}
The reductions are stated over $\Q$, where \cref{src:thm:dgrh} supplies the abelian input.
The same theorem gives the Riemann hypothesis for $\zeta_k$ for every finite abelian
extension $k/\Q$, and Lemma~\ref{nab:lem:heilbronn} holds over any base field. For the
detectors of $G=\PSL_2(\F_q)$, $q=2^r\ge4$, of \cref{grp:sec:psl-even} this gives a transfer over
abelian base fields. Let $M/\Q$ be Galois with group $G$ and $k/\Q$ finite abelian. Then
$M\cap k=\Q$, because its Galois group is an abelian quotient of the simple nonabelian
group $G$, so $\Gal(Mk/k)\simeq G$. The character identities and the nonnegativity of the
orders of the induced Hecke $L$-functions are unchanged with $k$ in place of $\Q$, the
logarithmic derivative now runs over prime ideals $\mathfrak p$ of $k$ and their powers
with coefficient $\log\Nm\mathfrak p$ times the trace, its principal part is still
$1/(s-1)$, and the Mellin calculation depends only on the norms. Since there are at most
$[k:\Q]$ prime ideals above each rational prime, the estimate
$O(\sqrt t\log^2(2t))$ for the omitted prime powers persists. Hence all the equivalences
of \cref{grp:sec:psl-even} for $M/\Q$ hold for $Mk/k$, and the one-sided criterion holds with
rational primes replaced by prime ideals of $k$, their norms replacing $p$, and root
counts of a defining polynomial of the relevant subfield $(Mk)^{U}=M^{U}k$ over $k$ taken
in the residue fields.
\end{remark}


\subsection{Symmetric and alternating groups: prime-order cycle patterns}\label{nab:sec:sn}
\begingroup\raggedright
\sources{grh_cyclic_generator_transfer_2026-11-25; grh_sylow_transfer_2026-11-25}
\endgroup

For $n\ge5$ one has $[S_n,S_n]=A_n=[A_n,A_n]$. Write $p^m1^{n-pm}$ for the cycle type with
$m$ disjoint $p$-cycles and all other points fixed. At a prime $p\notin S$ the cycle type
of $g_p$ is the list of degrees of the irreducible factors modulo $p$ of a polynomial of
degree $n$ with splitting field $E$.

\begin{theorem}[Sparse cycle patterns; \stU]\label{nab:thm:patterns}
Let $G=S_n$ or $A_n$, $n\ge5$, and $H=A_n$. The cyclic subgroups in
\eqref{nab:eq:weights-def} may be chosen with exactly one subgroup generated by an element
of each of the types
\begin{align*}
 &p^{p^j}1^{n-p^{j+1}}, && p\text{ an odd prime},\quad 0\le j<\floor{\log_pn};\\
 &2^{2^j}1^{n-2^{j+1}}, && 1\le j<\floor{\log_2n}.
\end{align*}
The support of $W^{\mathrm{cyc}}$ consists of these types and the identity. No Frobenius
element of composite order is needed, and there are
$\sum_{p\le n}\floor{\log_pn}-1=\pi(n)-1+O(\sqrt n)$ nonidentity patterns.
\end{theorem}
\begin{proof}
A Sylow $p$-subgroup of $S_{p^a}$ is the iterated wreath product
$Q_a=Q_{a-1}^p\rtimes C_p$, $Q_0=1$: its order $p^{1+p+\dots+p^{a-1}}$ is the $p$-part of
$(p^a)!$. The outer block rotation is a product of $p^{a-1}$ disjoint $p$-cycles, and
recursion in the blocks gives generators that are all rotations of this form at the
various levels. The base-$p$ expansion of $n$ gives disjoint blocks of sizes $p^a$, and the
product of the corresponding groups is a Sylow subgroup of $S_n$ by Legendre's formula.
For odd $p$ every generator is even, so this is also a Sylow subgroup of $A_n$.

For $p=2$ the Sylow subgroup of $S_n$ is $B\rtimes R$, where $B$ is the elementary abelian
group generated by the disjoint bottom transpositions $b_1,\dots,b_m$ and $R$ is generated
by the higher block rotations, which permute whole pairs without exchanging their two
labels and are therefore even. Its intersection with $A_n$ is $B_{\mathrm{even}}\rtimes R$,
and $B_{\mathrm{even}}$ is generated by the elements $b_1b_i$, $i>1$, of type $2^2$. Every
higher rotation has type $2^{2^j}$ with $j\ge1$. Since $n\ge5$, the type $2^2$ occurs, and
all listed levels occur in a largest block.

In $S_n$ the cyclic subgroups of prime order with a given generator type are conjugate.
The same holds in $A_n$ provided the normalizer in $S_n$ of such a subgroup contains an odd
permutation, since then every conjugation can be adjusted to an even one. For an
involution, one transposition of the generator centralizes it and is odd. For odd $p$ and
at least two $p$-cycles, the permutation exchanging two of the cycles by matching their
labels is a product of $p$ transpositions, hence odd, and centralizes the generator; if
there are two fixed points their transposition works as well. In the remaining case, one
$p$-cycle and at most one fixed point, identify the labels of the cycle with $\F_p$;
multiplication by a primitive root normalizes the subgroup and is a $(p-1)$-cycle on the
nonzero labels, hence odd. So one representative per type suffices for $A_n$ too.

For the count, the number of primes $p$ with $p^k\le n$ is at most $n^{1/k}$, and the sum
over $k\ge3$ is $O(n^{1/3}\log n)$.
\end{proof}

By \eqref{nab:eq:weights}, for $G=S_n$ the weight of $W^{\mathrm{cyc}}$ at a selected type
$p^m1^{n-pm}$ is
\[
 -(p-1)\,p^{m-1}\,m!\,(n-pm)! ,
\]
and for $G=A_n$ every weight, including that of the identity, is half of this; when a type
splits into two classes of $A_n$, half of the nonidentity powers lie in each class, so the
same formula applies. The integer weights may be divided by their greatest common divisor
without changing the criterion. For comparison, the support of $W^{\mathrm{Syl}}$ with
$H=A_n$ is the set of all nonidentity types of prime-power order in $A_n$. The numbers of
cycle types in the two supports, the identity excluded, are as follows
(\artifact{grh_cyclic_generator_transfer_2026-11-25/pattern_counts.csv}).
\begin{center}
\begin{tabular}{rrrrrrr}\toprule
$n$ & 8 & 10 & 16 & 20 & 30 & 40\\\midrule
cyclic-generator patterns & 5 & 6 & 9 & 11 & 15 & 18\\
Sylow patterns & 8 & 13 & 32 & 50 & 129 & 271\\\bottomrule
\end{tabular}
\end{center}

\begin{example}[Degrees four and five]\label{nab:ex:sylowsn}
For $G=S_4$ and $G=S_5$ with $H=A_n$, primitive integer multiples of
$W^{\mathrm{Syl}}$ are
\[
 17\cdot\nabone_{1^4}-3\cdot\nabone_{2^2}-\nabone_{1\,3},
 \qquad
 133\cdot\nabone_{1^5}-3\cdot\nabone_{1\,2^2}-2\cdot\nabone_{1^2\,3}-2\cdot\nabone_{5} .
\]
For $S_5$ the Sylow subgroups of $A_5$ are a Klein four-group, $C_3$ and $C_5$, of indices
$30,40,24$ in $S_5$, so the value at the identity is $360-94=266$; a double transposition
has centralizer of order $8$ and the four-group contains three of them, which gives
$-8\cdot3/4=-6$; a $3$-cycle and a $5$-cycle give $-6\cdot2/3=-4$ and $-5\cdot4/5=-4$; then
divide by $2$. For $S_4$ the indices are $6$ and $8$ and the nonidentity values are $-6$
and $-2$. The values are confirmed by direct enumeration in
\artifact{grh_sylow_transfer_2026-11-25/group_certificate.json}. The polynomial $X^5-X-1$
has discriminant $2869=19\cdot151$, is irreducible modulo $3$ and has type $(2,3)$ modulo
$2$. Its group is transitive of prime degree, hence primitive, and contains a
transposition, the cube of an element of type $(2,3)$; a primitive group containing a
transposition is the symmetric group. So this is an explicit instance with $G=S_5$, and
\cref{src:thm:dgrh} controls its quadratic subfield $\Q(\sqrt{2869})$. The corresponding
estimate \textup{(iii)} of \cref{nab:thm:sylow} is
\[
 \sum_{p\le x,\ p\nmid2869}\log p\Bigl(\frac1p-\frac1x\Bigr)
 \bigl(133\cdot\nabone_{1^5}-3\cdot\nabone_{1\,2^2}-2\cdot\nabone_{1^2\,3}-2\cdot\nabone_{5}\bigr)(p)
 \le c+C_\eps x^{-1/2+\eps},
\]
with the indicators referring to the factor degrees of the polynomial modulo $p$ and $c$
the centering constant. A worked instance of the cyclic-generator form in degree eight,
with six patterns, is \cref{grp:ex:s8} below.
\end{example}

\paragraph{A degree-eight target.}

\begin{example}[The splitting field of $T^8-T-1$; \stR]\label[example]{grp:ex:s8}
Let $f(T)=T^8-T-1$ with splitting field $E$. Modulo $7$, $19$ and $61$ the
factorization types are $(8)$, $(1,7)$ and $(1,2,5)$, with distinct factors;
the certificates are in\\
\artifact{grh_cyclic_generator_transfer_2026-11-25/certificate.json}.\\
Irreducibility modulo $7$ gives transitivity. The $7$-cycle gives primitivity:
in an invariant block system the block of its fixed point is stable under the
$7$-cycle, so has size $1$ or $8$. The fifth power of the element of type
$(1,2,5)$ is a transposition; the graph of its conjugates has components
forming a block system, hence is connected, and the transpositions along the
edges of a connected graph generate $S_8$. Therefore $\Gal(E/\Q)=S_8$.

The cyclic-generator form of \cref{nab:thm:sylow} uses the character
$W=\sum_CW_C$, the sum over one cyclic subgroup $C$ of each of the prime-order types
$2^2\,1^4$, $2^4$, $3\,1^5$, $5\,1^3$, $7\,1$ selected by \cref{nab:thm:patterns}. By
\eqref{nab:eq:weights}, with $W_C(1)=\abs G(1-\abs C^{-1})$, the weight $w=W/6$, as a function of the factorization pattern of $f$ modulo a prime
$\ell$ not dividing its discriminant, is
\begin{center}
\begin{tabular}{lrrrrrrr}\toprule
pattern & $1^8$ & $2^2\,1^4$ & $2^4$ & $3\,1^5$ & $5\,1^3$ & $7\,1$ & other\\\midrule
$w(\ell)$ & $22336$ & $-16$ & $-32$ & $-40$ & $-4$ & $-1$ & $0$\\\bottomrule
\end{tabular}
\end{center}
The class sizes are $1,210,105,112,1344,5760$, and the weighted sum vanishes.
With $c_w$ as in \cref{grp:prop:interface}, the inequality
\[
 \sum_{\ell\le x}w(\ell)\log\ell\Bigl(\frac1\ell-\frac1x\Bigr)-c_w
 \le C_\eps x^{-1/2+\eps}\qquad(\text{each }\eps>0)
\]
would, by \cref{nab:thm:sylow} and Dirichlet GRH for the quadratic subfield of
$E$, give the Riemann hypothesis for $\zeta_E$ and the absence of zeros and
poles of all its Artin factors off the critical line. It uses six patterns;
the Sylow-filtered weight for $S_8$ also involves $4\,2\,1^2$, $4^2$ and
$3^2\,1^2$.
\end{example}

\subsection{Reflection detectors}\label{nab:sec:reflection}
\begingroup\raggedright
\sources{grh_reflection_transfer_2026-11-25; grh_method_transfers_2026-11-26; grh_spanning_and_s5_2026-11-26}
\endgroup

For groups generated by reflections one subfield of index two in $E$ detects every zero.
Fix $\rho$ with $1/2<\Re\rho<1$, write $m=\ord_\rho\zeta_E$ and let $\Theta$ be the
Heilbronn character of $E/\Q$ at $\rho$. For a linear character $\epsilon:G\to\set{\pm1}$
the function $L(s,\epsilon,E/\Q)$ is a Dirichlet $L$-function, so \cref{src:thm:dgrh} gives
\begin{equation}\label{nab:eq:signorth}
 \langle\Theta,\epsilon\rangle=0 .
\end{equation}

\begin{theorem}[Local reflection inequality; \stG]\label{nab:thm:reflection}
Let $T\subseteq G$ be a conjugation-invariant set of involutions and $\epsilon$ a linear
character of $G$ with $\epsilon(t)=-1$ for $t\in T$. Suppose every $g\in G$ has an
expression $g=uv$ in which $u$ and $v$ are products of, respectively, $r_g$ and $s_g$
pairwise commuting members of $T$; empty products are allowed. Put
$a_t=\ord_\rho\zeta_{E^{\langle t\rangle}}$ for $t\in T$. For a chosen expression of each
$g$, with reflection list $t_{g,1},\dots,t_{g,r_g+s_g}$,
\begin{equation}\label{nab:eq:weighted}
 m\le\frac2{\abs G}\sum_{g\in G}\sum_{j=1}^{r_g+s_g}a_{t_{g,j}} .
\end{equation}
In particular, if every expression has at most $D$ reflections, then
$m\le2D\max_{t\in T}a_t$, and if all $a_t$ vanish then $m=0$.
\end{theorem}
\begin{proof}
By Lemma~\ref{nab:lem:heilbronn}, on an abelian subgroup containing $t$ the actual
representation with character $\Theta$ has a $+1$ eigenspace for $t$ of dimension
$a_t=(m+\Theta(t))/2$. For $u=t_1\cdots t_r$ with commuting factors the intersection of the
$-1$ eigenspaces of the $t_i$ has codimension at most $\sum_ia_{t_i}$. Hence the dimension
$b(u)$ of the space on which $u$ differs from the scalar $\epsilon(u)$ is at most
$\sum_ia_{t_i}$. This number equals $(m-\epsilon(u)\Theta(u))/2$ and depends only on the
character value, not on the chosen subgroup. The same applies to $v$.

The group $\langle u,v\rangle$ is dihedral, cyclic of order at most two, or trivial: $u$
and $v$ square to one and $u$ inverts $uv$. Take a unitary representation of this
subgroup with character $\Theta$. After twisting by $\epsilon$, the intersection of the
fixed spaces of $u$ and $v$ has codimension at most $b(u)+b(v)$, and $uv$ fixes this
intersection. All other eigenvalues of the unitary operator $uv$ have real part at least
$-1$. Therefore
\[
 m-\Re\bigl(\epsilon(g)\Theta(g)\bigr)\le2\bigl(b(u)+b(v)\bigr)\le2\sum_ja_{t_{g,j}} .
\]
Average over $G$; the average of $\epsilon(g)\Theta(g)$ is zero by
\eqref{nab:eq:signorth}. The argument compares character values on actual local
representations only and never assumes that $\Theta$ is positive definite.
\end{proof}

\begin{corollary}[Symmetric groups, all degrees; \stG]\label{nab:cor:sn}
Let $G=S_n$, $n\ge2$, $K=E^{\langle(12)\rangle}$ and $H_n=1+\tfrac12+\dots+\tfrac1n$. Then
$\ord_\rho\zeta_E\le2(n-H_n)\ord_\rho\zeta_K$. Hence, given \cref{src:thm:dgrh}, the
Riemann hypotheses for $\zeta_K$ and for $\zeta_E$ are equivalent.
\end{corollary}
\begin{proof}
For an $h$-cycle index the points by $\Z/h\Z$ with the cycle acting as $i\mapsto i+1$. Let
$u(i)=1-i$ and $v(i)=-i$; then $uv(i)=i+1$. For odd $h$ both involutions consist of
$(h-1)/2$ disjoint transpositions, and for even $h$ of $h/2$ and $h/2-1$. Factors on
distinct cycles commute, so $r_g+s_g=n-c(g)$ with $c(g)$ the number of cycles, fixed points
included. All transpositions are conjugate; take $\epsilon=\nabsgn$. A uniform random
permutation of $n$ letters has $H_n$ cycles on average: inserting the last letter into the
cycle notation of a permutation of $n-1$ letters, one of $n$ equally likely positions
creates a new singleton. Apply \eqref{nab:eq:weighted}. Conversely $E/K$ is quadratic and
its Hecke factor is holomorphic in the strip, so zeros of $\zeta_K$ are zeros of $\zeta_E$.
\end{proof}

\begin{corollary}[Finite real reflection groups; \stG]\label{nab:cor:coxeter}
Let $G$ be a finite real reflection group of rank $D$, acting in its reflection
representation. Choose one reflection $t_j$ from each conjugacy class of reflections and
put $K_j=E^{\langle t_j\rangle}$, $1\le j\le q$. Then
\[
 \ord_\rho\zeta_E\le2D\max_j\ord_\rho\zeta_{K_j}\le2D\sum_j\ord_\rho\zeta_{K_j} .
\]
Given \cref{src:thm:dgrh}, the Riemann hypothesis for $\zeta_E$ is equivalent to the
Riemann hypothesis for all $\zeta_{K_j}$. There is one class of reflections for the simply
laced irreducible Weyl groups and two for the types $B_n$ ($n\ge2$), $F_4$ and
$I_2(2k)$ ($k\ge2$).
\end{corollary}
\begin{proof}
By the factorization theorems of Carter and Springer \cite{Carter1972Conjugacy,Springer1974Regular}, in the length-additive form recorded in
\cite{BaumeisterDyerStumpWegener2017Hurwitz}, every $w\in G$ is a product $uv$ in which
each of $u,v$ is an involution or the identity and a product of commuting reflections,
with total number of reflections equal to the codimension of the fixed space of $w$, hence
at most $D$. The determinant is $-1$ on every reflection. Apply
\cref{nab:thm:reflection}. Each $E/K_j$ is quadratic, which gives the converse. The counts
of reflection classes follow from the conjugacy of simple reflections along paths with odd
labels in the Coxeter diagram.
\end{proof}

Without the length-additive statement, the factorization of each element into two
involutions and the decomposition of an involution into commuting reflections give at most
$2D$ factors, hence the same detection with constant $4D$. The symmetric-group case needs
neither the classification nor a finite computation.

The constant $2(n-H_n)$ is not optimal in small degrees. The next theorem gives the
constant $n-1$ for $3\le n\le12$ and shows that no argument of this kind can do better.

\begin{theorem}[Sharp multiplicities for symmetric groups of degree at most $12$; \stU, the bound over $\Q$ \stG]\label{nab:thm:snsharp}
Let $E/k$ be a Galois extension of number fields with group $S_n$, $3\le n\le12$. Fix a
transposition $t$, put $K=E^{\langle t\rangle}$ and $F=E^{A_n}$, and let $\varepsilon$ be
the sign character. Then
\[
 \mathcal Q(s)=\frac{\zeta_K(s)^{\,n-1}L(s,\varepsilon,E/k)}{\zeta_E(s)}
 =\frac{\zeta_K(s)^{\,n-1}\zeta_F(s)}{\zeta_E(s)\zeta_k(s)}
\]
is holomorphic on $\C\setminus\set1$ and has a pole of exact order $n-2$ at $s=1$. In
particular, for $k=\Q$ and $1/2<\Re\rho<1$, \cref{src:thm:dgrh} gives
\[
 \ord_\rho\zeta_E\le(n-1)\ord_\rho\zeta_K .
\]
For every $n\ge3$ the constant $n-1$ cannot be improved by any deduction that uses only
the nonnegativity of $\langle\Res_U\Theta,\eta\rangle$ for linear characters $\eta$ of
subgroups $U$ and the vanishing of the coefficients of the two linear characters of $S_n$.
\end{theorem}
\begin{proof}
Fix $s_0\ne1$, let $\Theta$ be the Heilbronn character of $E/k$ at $s_0$ and put
$\Psi=\varepsilon\Theta$. For every subgroup $U$ and every linear character $\eta$ of $U$,
Lemma~\ref{nab:lem:heilbronn} gives $\langle\Res_U\Theta,\eta\rangle\ge0$, and the same holds for
$\Psi$ because $\varepsilon|_U\eta$ is linear. Moreover
\[
 m:=\Psi(1)=\ord_{s_0}\zeta_E,\qquad
 a:=\tfrac12\bigl(m-\Psi(t)\bigr)=\ord_{s_0}\zeta_K,\qquad
 b:=\langle\Psi,\nabone\rangle=\ord_{s_0}L(s,\varepsilon,E/k).
\]
For a conjugacy class (cycle type) $T$ let $V_T$ be the value of $\Psi$ on $T$, $c_T$ the
size of $T$ and $\delta_T$ the corresponding unit vector. For a subgroup $U$ and a
surjection $e:U\to\Z/d$ put
\[
 B_{U,e}(T)=\sum_{h\in U\cap T}\frac1{\varphi(d)}\sum_{u\in(\Z/d)^\times}
   \exp\bigl(2\pi i\,u\,e(h)/d\bigr) .
\]
The row $B_{U,e}$ is real, since the sum over $u$ contains the complex conjugate of each
character, and $B_{U,e}\cdot V\ge0$ by the local positivity. For each $3\le n\le12$ there
are finitely many pairs $(U,e)$ and weights $w_{U,e}>0$ with
\begin{equation}\label{nab:eq:R}
 \frac{n-3}{n-1}\,\delta_1-\delta_t+\frac2{(n-1)\,n!}\,(c_T)_T=\sum_{(U,e)}w_{U,e}B_{U,e} .
\end{equation}
Pairing \eqref{nab:eq:R} with $V$ gives $2\bigl(a+(b-m)/(n-1)\bigr)\ge0$, that is
$m\le(n-1)a+b$ at every $s_0\ne1$, which is the holomorphy of $\mathcal Q$. At $s=1$ each
Dedekind zeta function has a simple pole and $L(1,\varepsilon,E/k)\ne0$, which gives the
exact order $n-2$. Over $\Q$ the factor $L(s,\varepsilon)$ is a Dirichlet $L$-function,
so $b=0$ for $1/2<\Re\rho<1$.

The identity \eqref{nab:eq:R} is a finite statement. For $n=3$ the single row is that of
the sign character of $S_3$, with weight $1/6$. For $n=5$ and $n=6$ the rows are given in
Table~\ref{nab:tab:rows}. For $n=5$, in the order of types $1,2,3,2^2,4,3\,2,5$ (fixed points
suppressed), with class sizes $(1,10,20,15,30,20,24)$, the four rows are
\[
\begin{pmatrix}
1&-2&0&1&0&0&0\\
1&-2&0&-1&2&0&0\\
1&-2&2&-3&0&2&0\\
1&0&0&-5&0&0&4
\end{pmatrix},
\]
and \eqref{nab:eq:R} reads
$\tfrac38r_1+\tfrac1{16}r_2+\tfrac1{24}r_3+\tfrac1{40}r_4
=\tfrac1{240}(121,-230,20,15,30,20,24)$. For the remaining degrees the pairs $(U,e)$, given
by generators with their images, and the weights are listed in
\artifact{grh_method_transfers_2026-11-26/checks/reflection/certificates.json}; the
program \artifact{grh_method_transfers_2026-11-26/checks/reflection/check_reflection.py}
enumerates each subgroup, sums the real cyclotomic traces in each conjugacy class and
checks \eqref{nab:eq:R} entry by entry in exact rational arithmetic. The numbers of rows
are in Table~\ref{nab:tab:rowcount}.

For the optimality, the standard character $\Psi(g)=\#\nabFix(g)-1$ of $S_n$, $n\ge3$, is
orthogonal to both linear characters, satisfies every subgroup positivity inequality, and
has $(m,a,b)=(n-1,1,0)$.
\end{proof}

\begin{table}[ht]
\centering\small
\begin{tabular}{rrrl}\toprule
$n$ & $w_{U,e}$ & $d$ & generators of $U$ with their images in $\Z/d$\\\midrule
5 & $3/8$ & 2 & $(45)\mapsto1,\ (23)\mapsto1$\\
5 & $1/16$ & 2 & $(45)\mapsto1,\ (24)(35)\mapsto1$\\
5 & $1/24$ & 2 & $(45)\mapsto1,\ (12)(34)\mapsto1$\\
5 & $1/40$ & 2 & $(23)(45)\mapsto1,\ (12)(34)\mapsto1$\\\midrule
6 & $7/30$ & 2 & $(56)\mapsto1,\ (23)(45)\mapsto1$\\
6 & $1/10$ & 3 & $(456)\mapsto2,\ (12)(34)(56)\mapsto0$\\
6 & $1/45$ & 6 & $(456)\mapsto4,\ (14)(25)(36)\mapsto3$\\
6 & $1/50$ & 2 & $(34)(56)\mapsto1,\ (23)(45)\mapsto1$\\
6 & $1/5$ & 2 & $(56)\mapsto1,\ (34)\mapsto1,\ (12)\mapsto1$\\
6 & $1/40$ & 2 & $(56)\mapsto0,\ (34)\mapsto1,\ (13)(24)\mapsto1$\\\bottomrule
\end{tabular}
\caption{Positive rows in \eqref{nab:eq:R} for $S_5$ and $S_6$.}\label{nab:tab:rows}
\end{table}

\begin{table}[ht]
\centering\small
\begin{tabular}{l|rrrrrrrrrr}
$n$&3&4&5&6&7&8&9&10&11&12\\\hline
conjugacy classes&3&5&7&11&15&22&30&42&56&77\\
positive subgroup rows&1&3&4&6&10&18&26&38&48&71
\end{tabular}
\caption{Sizes of the certificates for \cref{nab:thm:snsharp}.}\label{nab:tab:rowcount}
\end{table}

\begin{remark}[The character behind \cref{nab:thm:snsharp}]\label{nab:rem:un}
Let $\sigma$ be the nontrivial character of $\langle t\rangle$. The character whose
pairing with $\Psi$ is $(n-1)a-m+b$ is
\[
 U_n=(n-1)\Ind_{\langle t\rangle}^{S_n}\sigma-\Reg_{S_n}+\nabone .
\]
It is the character of an actual representation for every $n$. Indeed, in the Cayley
graph of $S_n$ for the $n-1$ adjacent transpositions, with $S_n$ acting by left
multiplication, the vertex module is regular and the module of oriented edges has one
summand $\Ind\sigma$ for each generator, because the stabilizer of an edge reverses it;
the graph is connected, so the cellular chain complex gives $U_n=H_1$ as characters. What
\cref{nab:thm:snsharp} needs is more: that $U_n$ lies in the positive rational cone spanned
by the characters induced from linear characters of subgroups. Formula \eqref{nab:eq:R}
supplies such a decomposition for $n\le12$, by Frobenius reciprocity. It does not follow
from graph homology alone. Let $A_5$ fix two vertices and permute five parallel edges in
its natural action; the first homology is the irreducible augmentation representation $U$
of dimension four. If $U$ were a positive rational combination of monomial characters,
each summand would be a multiple $kU$ with $1\le k\le4$ by Frobenius reciprocity, induced
from a subgroup of index $4k$. Indices $8$ and $16$ do not divide $60$, and a subgroup of
index $4$ would give a nontrivial homomorphism from the simple group $A_5$ to $S_4$. In
the remaining case the subgroup has order $5$ and $k=3$; but $U$ restricted to this
subgroup is its regular representation minus the trivial character, so every linear
character occurs at most once.
\end{remark}

The reflection inequalities turn into a test on rational primes with only two Frobenius
classes in the symmetric case.

\begin{theorem}[Reflection prime test; \stG; \stR]\label{nab:thm:reftest}
Let $G$ be a finite real reflection group, let $K_1,\dots,K_q$ be as in
Corollary~\ref{nab:cor:coxeter}, let $w_1,\dots,w_q>0$ be arbitrary fixed weights, and let
$R_{K_j}$ be the remainders of Proposition~\ref{nab:prop:dedekind}. Given \cref{src:thm:dgrh}, the
following are equivalent:
\begin{enumerate}[label=\textup{(\roman*)}]
\item the Riemann hypothesis for $\zeta_E$;
\item for every $\eta>0$ there is $C_\eta$ with
$\sum_{j=1}^qw_jR_{K_j}(x)\le C_\eta x^{-1/2+\eta}$ for all sufficiently large $x$.
\end{enumerate}
Only the single summed upper bound is required. For $G=S_n$ one has $q=1$, and the weight
$a_K(p)$ of the test is
\begin{equation}\label{nab:eq:snweight}
 a_K(p)=\begin{cases}n!/2,&g_p=1,\\ (n-2)!,&g_p\text{ a transposition},\\ 0,&\text{otherwise}.\end{cases}
\end{equation}
\end{theorem}
\begin{proof}
If $\zeta_E$ has a zero $\rho$ with $\Re\rho>1/2$, choose $0<\eta<\Re\rho-1/2$. By
Corollary~\ref{nab:cor:coxeter} some $\zeta_{K_j}$ vanishes at $\rho$, which contradicts
Proposition~\ref{nab:prop:dedekind}(a) with $\beta=1/2+\eta$. The functional equation gives (i).
Conversely the Riemann hypothesis for $\zeta_E$ implies it for every $\zeta_{K_j}$, and
Proposition~\ref{nab:prop:dedekind}(b) gives (ii). The number of fixed cosets of $g$ on
$G/\langle t\rangle$ is $\abs{\set{x\in G:x^{-1}gx\in\langle t\rangle}}/2$, which is
$\abs G/2$ at $g=1$, $\abs{C_G(t)}/2$ for $g$ conjugate to $t$, and zero otherwise; for a
transposition $\abs{C_{S_n}(t)}=2(n-2)!$. The mean of the weight is one, by Burnside's
lemma, in accordance with the main term $\log x$.
\end{proof}

In contrapositive form, a zero of $\zeta_E$ of real part $\sigma>1/2$ forces the left side
of (ii), divided by $x^{\beta-1}$, to be unbounded above for every $1/2<\beta<\sigma$.

\begin{example}\label{nab:ex:quintic}
For the splitting field $E$ of $f=X^5-X-1$ (Example~\ref{nab:ex:sylowsn}) one has
$[K:\Q]=60$, $[E:\Q]=120$, and the weight \eqref{nab:eq:snweight} is $60$ when $f$ splits
completely modulo $p$, $6$ when it has factor degrees $(1,1,1,2)$, and $0$ otherwise;
$60/120+6\cdot10/120=1$. A single upper bound for the centered smoothed sum of these
weights is equivalent, given \cref{src:thm:dgrh}, to the Riemann hypothesis for the field
of degree $120$. By \cref{nab:thm:snsharp} a zero of $\zeta_E$ of multiplicity $m$ off the
line has multiplicity at least $\ceil{m/4}$ in $\zeta_K$; Corollary~\ref{nab:cor:sn} gives
$\ceil{30m/163}$. For $S_6$ the two constants are $5$ and $71/10$.
\end{example}


\subsubsection{A three-subgroup identity in \texorpdfstring{$S_5$}{S5}}\label{grp:sec:s5}
Let $G=S_5$, $t=(12)$, let $\sigma$ be the nontrivial character of
$\langle t\rangle$, and $\grpstd$ the standard character,
$\grpstd(g)=\#\grpFix(g)-1$. Consider
\[
\begin{array}{c|c|c}
H&\text{generators}&\text{character}\\ \hline
C=C_6 &(12)(345)&\chi((12)(345))=e^{\pi i/3}\\
V=C_2\times C_2 &(12),\ (34)&\eta((12))=\eta((34))=-1\\
W=C_5\rtimes C_4 &a=(12345),\ b=(2354)&\xi(a)=1,\quad\xi(b)=i
\end{array}
\]
where $W$ is the affine group of $\F_5$, with $bab^{-1}=a^2$.

\begin{theorem}[Three-subgroup identity; \stU]\label{grp:thm:s5-identity}
\begin{equation}\label{grp:eq:s5-identity}
 \Ind_{\langle t\rangle}^{G}\sigma-\grpstd
 =\Ind_{C}^{G}\chi+\Ind_{V}^{G}\eta+\Ind_{W}^{G}\xi ,
\end{equation}
the three summands having degrees $20,30,6$. Both sides take the values
$(56,-8,-1,0,0,1,1)$ on the cycle types $1,\ 2,\ 3,\ 2^2,\ 4,\ 3\,2,\ 5$.
Let $\varepsilon$ be the sign character, $\tau=\varepsilon\,\grpstd$ and
$\alpha=(\varepsilon|_C)\chi$, of order three. Then
\begin{equation}\label{grp:eq:s5-twisted}
 \Ind_{\langle t\rangle}^{G}\triv-\tau
 =\Ind_{C}^{G}\alpha+\Ind_{V}^{G}\triv+\Ind_{W}^{G}\overline\xi .
\end{equation}
\end{theorem}

\begin{proof}
The sizes $c_T$ of the seven types, in the stated order, are
$1,10,20,15,30,20,24$. For a linear character $\theta$ of $H$ put
$B_{H,\theta}(T)=\sum_{h\in H\cap T}\theta(h)$. Direct enumeration gives
\[
\begin{array}{c|rrrrrrr}
 &1&2&3&2^2&4&3\,2&5\\ \hline
B_{C,\chi}&1&-1&-1&0&0&1&0\\
B_{V,\eta}&1&-2&0&1&0&0&0\\
B_{W,\xi}&1&0&0&-5&0&0&4
\end{array}
\]
For $C$ the two values on elements of order three sum to $-1$ and the two on
elements of order six to $1$. In $W$ the four nonidentity translations are
$5$-cycles of value $1$, the five elements with multiplier $-1$ are double
transpositions of value $-1$, and the ten elements of order four have five
values $i$ and five values $-i$. The induction formula at type $T$ is
$(\Ind_H^G\theta)(T)=\frac{120}{\abs Hc_T}B_{H,\theta}(T)$. Pairing the left
side of \eqref{grp:eq:s5-identity} against an arbitrary class function is the
functional with coordinate row
\[
 \frac12\delta_1-\frac12\delta_2
 -\Bigl(\frac{c_T(\#\grpFix(T)-1)}{120}\Bigr)_T
 =\Bigl(\frac7{15},-\frac23,-\frac16,0,0,\frac16,\frac15\Bigr),
\]
which equals $\frac16B_{C,\chi}+\frac14B_{V,\eta}+\frac1{20}B_{W,\xi}$ entry
by entry. The denominators are the subgroup orders, so this is
\eqref{grp:eq:s5-identity}. For \eqref{grp:eq:s5-twisted} multiply by
$\varepsilon$: on $V$ one has $\varepsilon\eta=\triv$, and on $W$, where $a$
is even and $b$ is odd, $\varepsilon\xi=\overline\xi$. The exact verification
is recorded in\\
\artifact{grh_spanning_and_s5_2026-11-26/checks/sn5_verification.json}.
\end{proof}

\begin{corollary}[Two entire quotients; \stU]\label[corollary]{grp:cor:s5-quotients}
Let $M/K$ be a Galois extension of number fields with group $S_5$, and put
\[
 E=M^{\langle t\rangle},\quad P=M^{S_4},\quad A=M^C,\quad F=M^V,\quad J=M^W,
\]
of degrees $60,5,20,30,6$ over $K$, with $S_4$ a point stabilizer. Then
\begin{equation}\label{grp:eq:s5-four-zeta}
 \mathcal Z(s)=\frac{\zeta_M(s)\zeta_K(s)}{\zeta_E(s)\zeta_P(s)}
 =L_A(s,\chi)L_F(s,\eta)L_J(s,\xi)
\end{equation}
is entire and nonzero at $s=1$, so that at every $\rho\in\C$
\begin{equation}\label{grp:eq:s5-order}
 \ord_\rho\zeta_E+\ord_\rho\zeta_P\le\ord_\rho\zeta_M+\ord_\rho\zeta_K .
\end{equation}
Moreover
\begin{equation}\label{grp:eq:s5-companion}
 \mathcal H(s)=\frac{\zeta_E(s)}{L_K(s,\tau,M/K)}
 =L_A(s,\alpha)\,\zeta_F(s)\,L_J(s,\overline\xi)
\end{equation}
is holomorphic on $\C\setminus\set1$ with a simple pole at $1$, and
\begin{equation}\label{grp:eq:s5-tau}
 \ord_\rho L_K(s,\tau,M/K)\le\ord_\rho\zeta_E\qquad(\rho\ne1).
\end{equation}
\end{corollary}

\begin{proof}
Induction of the regular representation of $C_2$ gives
$\Ind_{C_2}^G\sigma=\Reg_G-\Ind_{C_2}^G\triv$, and
$\grpstd=\Ind_{S_4}^G\triv-\triv$; so the left side of
\eqref{grp:eq:s5-identity} has Artin $L$-function $\mathcal Z$. Artin
induction and multiplicativity give \eqref{grp:eq:s5-four-zeta}, first for
$\Re s>1$ and then by continuation, with all local factors included. The
characters $\chi,\eta,\xi$ are nontrivial, of orders $6,2,4$, so the three
Hecke factors are entire and nonzero at $1$. The second statement follows in
the same way from \eqref{grp:eq:s5-twisted}; Brauer's theorem provides the
meromorphy of $L_K(s,\tau,M/K)$, the two nontrivial Hecke factors are nonzero
at $1$, and $\zeta_F$ has a simple pole there.
\end{proof}

Both statements are independent of the Riemann hypothesis and of Dirichlet
GRH. They do not prove that $L_K(s,\tau,M/K)$ is entire, since
\eqref{grp:eq:s5-tau} permits a negative left side, and they leave the zero
locations of $\zeta_E$ and $\zeta_M$ open. Four-zeta quotients for solvable
groups and normal subgroups occur in \cite{GunHazraSahu2025Holomorphy}; here
the subgroups are not normal.

\begin{remark}[The step in a general induction]\label[remark]{grp:rem:s5-induction}
For $n\ge4$ let $U_n$ be the character of \cref{nab:rem:un}, with $C_2=\langle t\rangle$,
and put
\[
 V_n=\Ind_{C_2}^{S_n}\sigma-\grpstd_n .
\]
Induction in stages gives $U_n=\Ind_{S_{n-1}}^{S_n}U_{n-1}+V_n$. The bound
$\ord_\rho\zeta_E\le(n-1)\ord_\rho\zeta_K$ of \cref{nab:thm:snsharp} rests on a positive
rational decomposition of $U_n$ into characters induced from linear characters, which
\eqref{nab:eq:R} certifies for $3\le n\le12$; a positive monomial decomposition of
every $V_n$ would extend it to all degrees. Equation
\eqref{grp:eq:s5-identity} is such a decomposition of $V_5$, with integral
coefficients.
\end{remark}

\subsection{Compression along normal towers and mixed factors of composita}\label{nab:sec:tower}
\begingroup\raggedright
\sources{grh_relative_prime_races_2026-11-25}
\endgroup

Let
\[
 \Q\subseteq F_0\subsetneq F_1\subsetneq\dots\subsetneq F_r=E,\qquad r\ge1,
\]
be a tower in which every field is Galois over $\Q$. Put $H_i=\Gal(E/F_i)$,
$d_i=[F_i:\Q]$ and $n_i=d_i/d_{i-1}\ge2$, so that $H_i\triangleleft G$ and $H_r=\set1$,
and let $Q_i=\zeta_{F_i}/\zeta_{F_{i-1}}$. By the Aramata--Brauer theorem
\cite{Brauer1947Zeta,MurtyMurty1997Nonvanishing} each $Q_i$ is entire; it is nonzero on
$\Re s\ge1$.

\begin{lemma}[Local coefficient]\label{nab:lem:local}
For $p\notin S$ the coefficient of $\log p\,p^{-ks}$ in $-\zeta_{F_i}'/\zeta_{F_i}$ is
$d_i\nabone_{H_i}(g_p^k)$. Consequently, for
$w_i(g)=n_i\nabone_{H_i}(g)-\nabone_{H_{i-1}}(g)$ the prime-power series
\[
  D_{w_i}(s)=\sum_{p\notin S}\sum_{k\ge1}w_i(g_p^k)\log p\,p^{-ks}
\]
of \cref{ana:thm:primeonly} satisfies $D_{w_i}(s)=-d_{i-1}^{-1}Q_i'(s)/Q_i(s)+H_i^*(s)$, where $H_i^*$ is a finite
combination of logarithmic derivatives of the Euler factors at the primes of $S$, holomorphic
and bounded on $\Re s\ge\eta$ for every $\eta>0$.
\end{lemma}
\begin{proof}
If the image of $g_p$ in $G/H_i$ has order $f$, then $p$ has $d_i/f$ prime ideals of norm
$p^f$ in $F_i$. Differentiating their Euler factors gives the coefficient $(d_i/f)f=d_i$
when $f\mid k$ and $0$ otherwise. The Euler factors at the primes of $S$ are finite
geometric series whose poles lie on $\Re s=0$.
\end{proof}

\begin{theorem}[Tower compression; \stR]\label{nab:thm:tower}
Put $c_1=1$ and $c_j=(n_1-1)\prod_{i=2}^{j-1}n_i$ for $2\le j\le r$. Then all $c_i$ are
positive and
\begin{equation}\label{nab:eq:compressed}
 w(g):=\sum_{i=1}^rc_iw_i(g)
 =(n_1-1)\frac{\abs{H_1}}{\abs{H_r}}\nabone_{H_r}(g)-\nabone_{H_0\setminus H_1}(g) .
\end{equation}
The relative Riemann hypothesis for $\zeta_E/\zeta_{F_0}$ is equivalent to the
constant-free inequality \textup{(v)} of \cref{ana:thm:primeonly} for this $w$, which has
mean zero: for every
$0<\eps<1/2$ and all sufficiently large integers $x$,
$P_w(x)-P_w(x^2)\le C_\eps x^{-1/2+\eps}$. If $F_0/\Q$ is abelian, \cref{src:thm:dgrh}
makes this a criterion for the Riemann hypothesis of $\zeta_E$ itself.
\end{theorem}
\begin{proof}
On $H_0\setminus H_1$ the sum is $-1$. On $H_{j-1}\setminus H_j$ with $j\ge2$ it is
$\sum_{i<j}c_i(n_i-1)-c_j=0$, because the recursion defining the $c_j$ is exactly
$c_j=\sum_{i<j}c_i(n_i-1)$. On $H_r$ it is $\sum_{i\le r}c_i(n_i-1)=(n_1-1)\prod_{i=2}^rn_i$,
and outside $H_0$ all terms vanish. By Lemma~\ref{nab:lem:local},
$D_w=-\sum_i(c_i/d_{i-1})Q_i'/Q_i+H$ with positive coefficients, and each $w_i$ has mean
$n_i/d_i-1/d_{i-1}=0$, so \cref{ana:thm:primeonly}\textup{(i)}$\Leftrightarrow$\textup{(v)}
applies with $a_w=0$. The entire quotients $Q_i$ have product
$\zeta_E/\zeta_{F_0}$ and cannot cancel one another's zeros; their zeros in the strip are
symmetric under the functional equations. The last assertion follows from the
factorization of $\zeta_{F_0}$ into Dirichlet $L$-functions.
\end{proof}

\begin{remark}[Support]\label{nab:rem:towersupport}
The compressed weight sees only the primes that split completely in $E$ and the primes
that split completely in $F_0$ but not in $F_1$. Its support has density
$(1+\abs{H_0}-\abs{H_1})/\abs G$, against $\abs{H_0}/\abs G$ for the direct quotient race
$([E:F_0]-1)\nabone_{\set1}-\nabone_{H_0\setminus\set1}$; the compression removes
$\abs{H_1}-1$ group elements from the support, at the price of a large weight at the
identity. Among strictly positive combinations of the $w_i$, the first stratum and the
identity must remain, and the cancellation of every other stratum determines the ratios
$c_i/c_1$. This is an optimality statement inside that cone only. For $S_4$ with the tower
$\Q(\sqrt{\disc})\subset E^{V_4}\subset E$ the compressed weight is
$8\cdot\nabone_{1^4}-\nabone_{1\,3}$; see \cref{grp:thm:s4-race}.
\end{remark}

\begin{proposition}[Mixed factors of composita; \stR]\label{nab:prop:composita}
Let $E_1,E_2$ be Galois over $\Q$, both containing a field $F$ Galois over $\Q$, with
$E_1\cap E_2=F$. Put $E=E_1E_2$, $d=[F:\Q]$ and $n_i=[E_i:F]>1$. The function
\[
 Q_{12}(s)=\frac{\zeta_E(s)\zeta_F(s)}{\zeta_{E_1}(s)\zeta_{E_2}(s)}
\]
is entire by a theorem of Brauer \cite{Brauer1973Note}, and its Riemann hypothesis is
equivalent to the constant-free inequality \textup{(v)} of \cref{ana:thm:primeonly} for the
class function of mean zero
on $\Gal(E/\Q)$
\[
 w_{12}(g)=\nabone_{g|_F=1}\bigl(n_1\nabone_{g|_{E_1}=1}-1\bigr)\bigl(n_2\nabone_{g|_{E_2}=1}-1\bigr).
\]
The same holds for more than two fields that are jointly linearly disjoint over $F$, with
the product of the factors $n_i\nabone_{g|_{E_i}=1}-1$.
\end{proposition}
\begin{proof}
Expanding the product gives
$n_1n_2\nabone_{g=1}-n_1\nabone_{g|_{E_1}=1}-n_2\nabone_{g|_{E_2}=1}+\nabone_{g|_F=1}$,
and the computation of Lemma~\ref{nab:lem:local} gives $D_{w_{12}}=-Q_{12}'/(dQ_{12})+H$.
$Q_{12}$ is nonzero on $\Re s\ge1$ and has a functional equation by multiplication. Apply
\cref{ana:thm:primeonly}\textup{(i)}$\Leftrightarrow$\textup{(v)} with $a_w=0$. For several fields the entireness is Brauer's
inclusion--exclusion theorem; the hypothesis is joint disjointness, not pairwise
disjointness.
\end{proof}

The criterion does not require the Riemann hypothesis for either constituent field. For
$F=\Q$ and $G=G_1\times G_2$ the weight is the character of the tensor product of the two
augmentation representations, and the test detects the mixed zeros without using the
holomorphy of the individual tensor-product Artin $L$-functions. Even if both $\zeta_{E_1}$
and $\zeta_{E_2}$ satisfy the Riemann hypothesis, the mixed factor is not controlled by
that alone.

\subsection{Nonnegative deficits and the nonsplit-prime test}\label{nab:sec:deficit}
\begingroup\raggedright
\sources{grh_deficit_transfer_2026-11-25; grh_sparse_singular_moduli; grh_sharp_torsion_and_symplectic_blocks; mcnew_density_and_grh_transfers}
\endgroup

The geometric part of the source proof attaches nonnegative local amounts to primes
(\cref{src:sec:mechanism}). The next result gives, for a prescribed combination of
permutation characters, the least nonnegative Frobenius weight obtained from it by an
abelian correction, and its one-sided test.

Let $H=[G,G]$ and $B=E^H$; \cref{src:thm:dgrh} applies to $\zeta_B$ and to all characters
of $G/H$. Let $K_j=E^{U_j}$ be subfields, with no normality assumption on $U_j$, let
$\pi_j=\Ind_{U_j}^G\nabone$, fix $b_j>0$ and put
\[
 \phi(g)=\sum_jb_j\pi_j(g),\qquad M(q)=\max_{gH=q}\phi(g),\qquad w(g)=M(gH)-\phi(g).
\]
Then $w$ is a nonnegative real class function, and it is the pointwise least nonnegative
function of the form $f(gH)-\phi(g)$, since nonnegativity forces $f(q)\ge M(q)$ on every
coset $q$. In particular it minimizes every sum $\sum_gc(g)w(g)$ with $c(g)\ge0$, and its
support is contained in the support of every other such function.

\begin{theorem}[Permutation-deficit criterion; \stG; \stR]\label{nab:thm:deficit}
With $w$ as above, fix $c>1$, let $K_c$ be the kernel of \cref{nab:sec:interface}, and
put
\[
 \delta=\frac1{\abs G}\sum_{g\in G}w(g),\qquad
 D(x)=\frac1x\sum_{\substack{p\le cx\\p\notin S}}w(g_p)\log p\,K_c(p/x)-\delta\log c .
\]
Given \cref{src:thm:dgrh}, the following are equivalent:
\begin{enumerate}[label=\textup{(\roman*)}]
\item the Riemann hypothesis for every $\zeta_{K_j}$;
\item $D(x)=O_\eps(x^{-1/2+\eps})$ for every $\eps>0$;
\item $D(x)\le C_\eps x^{-1/2+\eps}$ for every $\eps>0$ and all sufficiently large $x$;
\item $\int_X^{2X}\nabpos{D(x)}\,\dd x=O_\eps(X^{1/2+\eps})$ for every $\eps>0$.
\end{enumerate}
\end{theorem}
\begin{proof}
Expand $M(q)=\sum_\chi m_\chi\chi(q)$ in characters of the abelian group $G/H$. The Euler
coefficients of $-\zeta_{K_j}'/\zeta_{K_j}$ at $p^k$, $p\notin S$, are
$\pi_j(g_p^k)\log p$, since the cycle lengths of Frobenius on $G/U_j$ are the residue
degrees. By the Kronecker--Weber theorem and Artin reciprocity for $G/H$, the series
$A(s)=\sum_{p\notin S,k\ge1}w(g_p^k)\log p\,p^{-ks}$ equals, up to finite Euler corrections
holomorphic in $\Re s>0$,
\[
 \sum_{\chi\in\widehat{G/H}}m_\chi\Bigl(-\frac{L'(s,\chi)}{L(s,\chi)}\Bigr)
 +\sum_jb_j\frac{\zeta_{K_j}'(s)}{\zeta_{K_j}(s)} .
\]
The mean of each transitive permutation character is one, so the residue at $1$ is
$m_{\nabone}-\sum_jb_j=\delta$. By \cref{src:thm:dgrh} the first sum has no other pole in
$\Re s>1/2$. At a zero $\rho$ of some $\zeta_{K_j}$ in that region the residue of the
second sum is $\sum_jb_j\ord_\rho\zeta_{K_j}>0$, without cancellation. Hence the
coefficients $-a$ of $-A$ satisfy the hypotheses of \cref{ana:thm:fd} at $\beta=1/2$, and the
statistic $\Delta_c$ of $-a$ is the prime-power version of $D(x)$; by \cref{ana:lem:powers}
the two differ by $O(x^{-1/2})$. So \cref{ana:thm:fd}, in the orientation described after it,
shows that (iv) implies (i), using the functional equations. Conversely, under (i) all the $L$-functions involved have their zeros on the
critical line, and shifting the contour in
$T(cx)-T(x)=(2\pi i)^{-1}\int_{(2)}A(s)H_c(s)x^{s-1}\,\dd s$ to $\Re s=1/2+\eps$, with the
classical bound $O_\eps(\log^2(3+\abs t))$ for the logarithmic derivatives, gives (ii); see
\cref{sec:ana}. The implications (ii)$\Rightarrow$(iii)$\Rightarrow$(iv) are immediate. No
Artin holomorphy is used.
\end{proof}

The regular action normalized by $\abs G$, that is $\phi=\nabone_{\set1}$, gives
$w=\nabone_{H\setminus\set1}$: a single unweighted test on the primes that split completely
in $B$ but not in $E$, in which completely split primes do not occur. This is form (v) of
the following theorem, which collects the tests of this event, with other smoothings and
for any normal subgroup in place of $[G,G]$.

\begin{theorem}[The nonsplit-prime test; \stR]\label{nab:thm:nonsplit}
Let $N\ne1$ be a normal subgroup of $G$, $F=E^N$, $a=[F:\Q]$ and $g=\abs G$. Assume that
$\zeta_F$ has no zeros in $\Re s>1/2$; by \cref{src:thm:dgrh} this holds when $G/N$ is
abelian, in particular for $N=[G,G]$. Let $\mathcal P$ be the set of primes $p\notin S$
that split completely in $F$ and not in $E$, that is, with $g_p\in N\setminus\set1$; its
density is $\delta=1/a-1/g$. Put
\[
 R(x)=\sum_{\substack{p\le x\\p\in\mathcal P}}\log p\Bigl(\frac1p-\frac1x\Bigr)-\delta\log x-c,
\]
where $c$ is the constant for which $R(x)\to0$. The following are equivalent:
\begin{enumerate}[label=\textup{(\roman*)}]
\item the Riemann hypothesis for $\zeta_E$;
\item there is $C$ with $R(x)\ge-Cx^{-1/2}$ for all sufficiently large $x$;
\item for every $0<\eps<1/2$ there is $C_\eps$ with $R(x)\ge-C_\eps x^{-1/2+\eps}$ for all
sufficiently large $x$.
\end{enumerate}
When $N=[G,G]$ these are also equivalent, given \cref{src:thm:dgrh}, to each of:
\begin{enumerate}[label=\textup{(\roman*)},resume]
\item for every $\eps>0$,
$\sum_{p\le x,\,p\in\mathcal P}\log p-\delta x\le C_\eps x^{1/2+\eps}$;
\item the finite-difference statistic $D(x)$ of \cref{nab:thm:deficit} for
$w=\nabone_{N\setminus\set1}$ satisfies $D(x)\le C_\eps x^{-1/2+\eps}$ for every $\eps>0$
and all sufficiently large $x$.
\end{enumerate}
\end{theorem}
\begin{proof}
Since $F$ and $E$ are Galois over $\Q$, $a_F(p)=a\nabone_N(g_p)$ and
$a_E(p)=g\nabone_{\set1}(g_p)$, so $\nabone_{\mathcal P}(p)=a_F(p)/a-a_E(p)/g$ and, in the
notation of Proposition~\ref{nab:prop:dedekind},
\[
 R(x)=\frac1aR_F(x)-\frac1gR_E(x),\qquad c=\frac{c_F}a-\frac{c_E}g .
\]
Equivalently, the series $B(s)=\sum_{p\in\mathcal P}\log p\,p^{-s}$ equals
$-a^{-1}\zeta_F'/\zeta_F+g^{-1}\zeta_E'/\zeta_E$ up to a function holomorphic in
$\Re s>1/2$, and the Mellin transform $B(s)/(s(s-1))-\delta/(s-1)^2-c/(s-1)$ of $R$ has
residue $\ord_\rho\zeta_E/(g\rho(\rho-1))$ at a zero $\rho$ of $\zeta_E$ in
$1/2<\Re s<1$, which is negative for real $\rho$. By hypothesis and
Proposition~\ref{nab:prop:dedekind}(b), $R_F(x)=O(x^{-1/2})$. If (i) holds, then also
$R_E(x)=O(x^{-1/2})$, which gives (ii), and (ii) implies (iii) trivially. If (iii) holds,
then $R_E(x)\le g\,(R_F(x)/a+C_\eps x^{-1/2+\eps})\le C_\eps'x^{-1/2+\eps}$, and
Proposition~\ref{nab:prop:dedekind}(a) excludes the zeros of $\zeta_E$ in $\Re s>1/2+\eps$; let
$\eps\to0$ and use the functional equation.

Now let $N=[G,G]$. Statement (v) is equivalent to (i) by \cref{nab:thm:deficit} with
$\phi=\nabone_{\set1}$. For (iv), put $f=\nabone_N-\nabone_{\set1}$ and
$\Psi(x)=\sum_{p^j\le x,\,p\notin S}f(g_p^j)\log p$; the terms with $j\ge2$ contribute
$O(\sqrt x\log^2(2x))$. For $\Re s>1$,
\[
 \int_1^\infty\bigl(\Psi(x)-\delta x\bigr)x^{-s-1}\,\dd x
 =\frac1s\Bigl(-\frac{\zeta_F'(s)}{a\,\zeta_F(s)}+\frac{\zeta_E'(s)}{g\,\zeta_E(s)}+U(s)\Bigr)
  -\frac{\delta}{s-1},
\]
with $U$ holomorphic in $\Re s>0$, because the normalized coefficient of
$-\zeta_M'/\zeta_M$ at $p^j$, divided by $\log p$, is one exactly when $g_p^j$ is trivial on
the Galois field $M$. The pole at $1$ cancels, $\zeta_F$ has no zero in $\Re s>1/2$, and at
a zero $\rho$ of $\zeta_E$ there the residue is $\ord_\rho\zeta_E/(g\rho)$, positive when
$\rho$ is real. Given (iv), fix $\theta=1/2+\eps<1$; then $Cx^\theta-\Psi(x)+\delta x\ge0$
for a suitable $C$, and \cref{ana:thm:onesided} shows that its transform is holomorphic in
$\Re s>\theta$, which excludes the zeros of $\zeta_E$ there. Conversely, under the Riemann
hypothesis for $\zeta_F$ and $\zeta_E$ one has
$\psi_M(x)=x+O_M(\sqrt x\log^2(2x))$ for $M=F,E$ \cite{GrenieMolteni2016Explicit};
subtracting the normalized counts gives (iv), the omitted primes contributing
$O(\log(2x))$.
\end{proof}

In (ii)--(iii) the test is a lower bound and in (iv)--(v) an upper bound, in accordance
with the signs of the kernels noted in \cref{nab:sec:interface}. For normal subgroups
$K\subsetneq N$ of $G$, the theorem applied to the Galois field $E^K$, with group $G/K$ and
normal subgroup $N/K$, shows that the event ``$p$ splits completely in $F$ and not in
$E^K$'', that is $g_p\in N\setminus K$, tests exactly the Riemann hypothesis for
$\zeta_{E^K}$: the ordinary Frobenius test has no access to data distinguishing the lifts
to $G$ of an element of $G/K$. This is the arithmetic meaning of the support-budget theorem
in \cref{nab:thm:budget}.

\begin{example}\label{nab:ex:nonsplit}
(a) Let $E=\Q(i,\sqrt[4]2)$, the splitting field of $X^4-2$, with dihedral group of order
$8$. The commutator subgroup has order $2$ and fixes $F=\Q(\zeta_8)$. For odd $p$ the event
$p\in\mathcal P$ is
\[
 p\equiv1\pmod 8,\qquad 2^{(p-1)/4}\equiv-1\pmod p,
\]
since for $p\equiv1\pmod8$ the number $2$ is a square and all fourth roots of unity are
available, so that failure of complete splitting in $E$ is the quartic nonresidue
condition. The first such primes are $17,41,97,137,193$ and $\delta=1/8$. A lower estimate
for the smoothed count with main term $(\log x)/8$ and error $O(x^{-1/2+\eps})$ is
equivalent to the Riemann hypothesis for $\zeta_E$.

(b) Let $E$ be the splitting field of $X^4-X-1$, of discriminant $-283$ and group $S_4$,
with $F=\Q(\sqrt{-283})$ and $M=E^{V_4}$ the splitting field of the cubic resolvent
$Y^3+4Y-1$. For $p\ne283$, the primes with factor degrees $(3,1)$, of density $1/3$, test
the Riemann hypothesis for $\zeta_M$; the primes with factor degrees $(3,1)$ or $(2,2)$, of
density $11/24$, test it for $\zeta_E$. The first prime distinguishing the two events is
$71$, of type $(2,2)$; among the $167$ unramified primes up to $1000$ the first event
selects $59$ and the second $82$. The counts are reproduced by
\artifact{grh_sharp_torsion_and_symplectic_blocks/checks/check_quartic.py}. The second
event, unweighted and tested by a finite difference, is the unweighted case of
\cref{grp:thm:s4}, for the same field.
\end{example}

\subsection{The zeros missing from the abelian part of a number field}\label{nab:sec:missing}
\begingroup\raggedright
\sources{grh-arithmetic-transfers; grh_dihedral_47_71}
\endgroup

For a number field $K$ that is not normal, the abelian part of its zeta function can be
removed in the same way, with a kernel that needs no centering constant. Let $E$ be the
normal closure of $K$, $U=\Gal(E/K)$, and $F=E^{U[G,G]}$, the largest subfield of $K$ that
is abelian over $\Q$. Put
\[
 U_K(x)=\sum_{\substack{p\le x\\p\notin S}}\bigl(a_K(p)-a_F(p)\bigr)\log p\,\Bigl(1-\frac px\Bigr).
\]

\begin{theorem}[One-sided criterion for the missing zeros; \stG; \stR]\label{nab:thm:missing}
Given \cref{src:thm:dgrh},
\[
 \mathrm{RH}(\zeta_K)\quad\Longleftrightarrow\quad
 \exists\,C_K\ \ \forall n\in\Z,\ n\ge2:\quad U_K(n)\ge-C_K\sqrt n .
\]
\end{theorem}
\begin{proof}
For a field $J$ let
$V_J(x)=\sum_{\Nm\mathfrak p^r\le x}\log\Nm\mathfrak p\,(1-\Nm\mathfrak p^r/x)
=x^{-1}\int_1^x\psi_J(u)\,\dd u$. Under the Riemann hypothesis for $\zeta_J$ the integrated
explicit formula gives $V_J(x)=x/2+O_J(\sqrt x)$, the sum over zeros being bounded by
$\sqrt x\sum_\rho\abs{\rho(\rho+1)}^{-1}$ \cite{GrenieMolteni2016Smoothed}. Prime powers
with rational exponent at least two contribute $O_J(\sqrt x)$, by the Chebyshev bound for
squares and then for higher exponents, and the primes of $S$ cost $O_J(\log x)$. Thus
$U_K=V_K-V_F+O(\sqrt x)$ unconditionally, and since $\zeta_F$ satisfies the Riemann
hypothesis by its factorization into Dirichlet $L$-functions, $\mathrm{RH}(\zeta_K)$ gives
$U_K(x)=O(\sqrt x)$.

Conversely, extend the lower bound from integers to all real $x\ge2$: $U_K$ is continuous
and, between primes, $\abs{U_K'(x)}\ll_Kx^{-2}\sum_{p\le x}p\log p\ll_K1$. Write
$W=V_K-V_F$. For $\Re s>1$,
\[
 \int_1^\infty W(x)x^{-s-1}\,\dd x
 =\frac{-\zeta_K'(s)/\zeta_K(s)+\zeta_F'(s)/\zeta_F(s)}{s(s+1)},
\]
and the poles at $s=1$ cancel. A zero $\rho$ of $\zeta_K$ with $\Re\rho>1/2$ gives a pole,
because $\zeta_F$ has no zero there; at a real zero $\beta\in(1/2,1)$ of multiplicity $h$
the residue is $-h/(\beta(\beta+1))<0$. The function $W(x)+C\sqrt x$ is nonnegative for a
suitable $C$, and its transform differs from that of $W$ by $C/(s-1/2)$. By
\cref{ana:thm:onesided} the transform of a nonnegative function cannot have a pole of
negative residue at its real abscissa of convergence, and it is regular at $1$; hence it is
holomorphic in $\Re s>1/2$ and $\zeta_K$ has no zero there. The functional equation gives
the other half of the strip.
\end{proof}


The weight $a_F$ subtracted here is the best approximation to $a_K$ by a function of a
residue class of $p$ (Corollary~\ref{nab:cor:doublecoset}). For a primitive sextic field with
normal closure of group $A_5$ one has $F=\Q$, and the weights $a_K(p)-1$ are $5,1,-1,0$ at
Frobenius elements of orders $1,2,3,5$. For a cubic field with group $S_3$ the weight is
the trace $\chi_\tau=(2,0,-1)$ of the standard representation $\tau$. In this case the
criterion can be rewritten with the least nonnegative weight
$w_*=\nabone_{A_3\setminus\set1}$ of Proposition~\ref{nab:prop:least}, for which
$3w_*=1+\nabsgn-\chi_\tau$: it becomes an upper bound for the Riesz mean of the primes
modulo which a defining cubic is irreducible, the Riesz-mean form of
\cref{nab:thm:nonsplit} for $N=A_3$. That criterion, for a cubic polynomial $P$ as in
\cref{hec:thm:irred}, is \cref{hec:prop:irredriesz}.


\subsection{Which Frobenius weights can work: support, sparsity and signs}\label{nab:sec:extremal}
\begingroup\raggedright
\sources{grh_sparse_singular_moduli; grh_dihedral_47_71; grh_sharp_torsion_and_symplectic_blocks; mcnew_density_and_grh_transfers; mcnew_counterexample_and_grh_transfers}
\endgroup

This subsection is finite group theory; no analytic hypothesis enters. It answers the
following question: once a quotient $G/N$ has been controlled, how few Frobenius classes
can a test see while every remaining irreducible character enters with the sign required
by the one-sided method? The results concern tests designed by the signs of individual
character coefficients. Arithmetic relations among zero multiplicities, such as those of
\cref{nab:sec:heilbronn,nab:sec:reflection}, allow other tests, and signed weights such as
those of \cref{nab:thm:sylow,nab:thm:tower} are not covered.

Let $G$ be a finite group of order $g$ and $N\ne1$ a normal subgroup of order $n$. An
irreducible complex character $\chi$ is called \emph{new} if its representation is
nontrivial on $N$; the other irreducible characters are inflated from $G/N$. For
$N=[G,G]$ the new characters are the nonlinear ones. Let $f$ be a real class function with
$f(u^{-1})=f(u)$, and write $f=\sum_\chi a_\chi\chi$ with $a_\chi=\langle f,\chi\rangle$;
the coefficients are real, and conjugate characters have equal coefficients.

\begin{lemma}\label{nab:lem:identity}
For $u\in N$,
\begin{equation}\label{nab:eq:identity}
 f(u)-f(1)=\sum_{\chi\ \mathrm{new}}(-a_\chi)\bigl(\chi(1)-\Re\chi(u)\bigr),
\end{equation}
and each bracket is nonnegative, with equality exactly when $u\in\ker\chi$. Moreover, for
new $\chi$ the mean of $\chi$ over $N$ is zero, $\sum_{\chi\ \mathrm{new}}\chi(1)^2=g-g/n$,
and
\begin{equation}\label{nab:eq:regular}
 \frac1g\sum_{\chi\ \mathrm{new}}\chi(1)\bigl(\chi(1)-\Re\chi(u)\bigr)=1-\nabone_{\set1}(u)
 \qquad(u\in N).
\end{equation}
\end{lemma}
\begin{proof}
The characters inflated from $G/N$ take their full degree on $N$ and cancel in
$f(u)-f(1)$; taking real parts gives \eqref{nab:eq:identity}. A unitary matrix of size
$\chi(1)$ has real trace at most $\chi(1)$, with equality only for the identity matrix.
Averaging a representation over $N$ projects onto its $N$-invariants, a $G$-submodule,
which is zero for a new irreducible representation. The remaining statements follow from
the regular character, $\sum_\chi\chi(1)\chi(u)=g\nabone_{\set1}(u)$, in which the
characters of $G/N$ contribute $g/n$ at every $u\in N$.
\end{proof}

\begin{theorem}[Exact support and minimax; \stU]\label{nab:thm:minimax}
Suppose $0\le f\le1$, $f(1)=0$ and $a_\chi\le0$ for every new $\chi$. Put
$t(f)=\min_{\chi\ \mathrm{new}}(-a_\chi)/\chi(1)$.
\begin{enumerate}[label=\textup{(\alph*)}]
\item The zero set of $f$ on $N$ is $K=N\cap\bigcap_{\chi\ \mathrm{new},\,a_\chi<0}\ker\chi$.
It is a normal subgroup of $G$, and $f$ factors through $G/K$.
\item If every new coefficient is strictly negative, the support of $f$ contains
$N\setminus\set1$. The indicator $f_*=\nabone_N-\nabone_{\set1}$ attains this minimal
support.
\item $\max_ft(f)=1/g$. The maximizers are exactly the functions $f_*+h$ with $h$ a real
inversion-symmetric class function inflated from $G/N$, vanishing on $N$, and with
$0\le f_*+h\le1$. Thus $f_*$ is the unique maximizer of minimal support.
\item If instead $f(1)=0$, $f\ge0$ and $a_\chi\ge0$ for every new $\chi$, then all new
coefficients vanish.
\end{enumerate}
\end{theorem}
\begin{proof}
(a) By \eqref{nab:eq:identity} with $f(1)=0$, $f(u)=0$ for $u\in N$ exactly when $u$ lies
in the kernel of every new $\chi$ with $a_\chi<0$. Every character occurring in $f$ is
trivial on $K$: the active new ones by definition and the others because they are inflated
from $G/N$. (b) The regular representation is faithful, so the new irreducible characters
together separate the identity from every other element of $N$; and
$a_\chi(f_*)=-\chi(1)/g$ for new $\chi$ because $\nabone_N$ has only constituents inflated
from $G/N$ and $\nabone_{\set1}=\Reg_G/g$. (c) By Lemma~\ref{nab:lem:identity},
\[
 t(f)\Bigl(g-\frac gn\Bigr)\le\sum_{\chi\ \mathrm{new}}(-a_\chi)\chi(1)
 =\frac1n\sum_{u\in N}f(u)\le1-\frac1n .
\]
Both estimates are sharp for $f_*$. Equality in the first forces $a_\chi=-\chi(1)/g$ for
every new $\chi$; then $f-f_*$ is inflated from the quotient, and $f(1)=f_*(1)=0$ forces it
to vanish on $N$. The converse is the same computation. (d) Average
\eqref{nab:eq:identity} over $N$: the left side is nonnegative and the right side is
$-\sum_{\chi\ \mathrm{new}}a_\chi\chi(1)\le0$.
\end{proof}

In the normalization natural for a nonnegative weight that is not bounded above, the same
identities give a pointwise statement.

\begin{proposition}[The least nonnegative weight; \stU]\label{nab:prop:least}
Let $N=[G,G]\ne1$, write $d_\pi=\chi_\pi(1)$, and let $w=\sum_\pi a_\pi\chi_\pi\ge0$ be a
real inversion-symmetric class function with $-a_\pi\ge d_\pi/\abs G$ for every nonlinear
$\pi$. Then $w\ge w_*:=\nabone_{N\setminus\set1}$ pointwise. More precisely, with
$b_\pi=-a_\pi-d_\pi/\abs G\ge0$, for $h\in N$
\[
 w(h)-w(1)-w_*(h)=\sum_{d_\pi>1}b_\pi\bigl(d_\pi-\Re\chi_\pi(h)\bigr),\qquad
 \frac1{\abs N}\sum_{h\in N}\bigl(w(h)-w(1)-w_*(h)\bigr)=\sum_{d_\pi>1}d_\pi b_\pi .
\]
Equality $w=w_*$ is attained, and $w_*$ is the unique minimizer of every objective
$\sum_gc(g)w(g)$ with all $c(g)>0$.
\end{proposition}
\begin{proof}
Subtract \eqref{nab:eq:regular} from \eqref{nab:eq:identity} to get the first identity;
with $w(1)\ge0$ and $w\ge0$ outside $N$ it gives $w\ge w_*$. Averaging over $N$ and using
the vanishing of the mean of a nonlinear character on $N$ gives the second. Attainment is
$a_\pi(w_*)=-d_\pi/\abs G$, as in \cref{nab:thm:minimax}(b).
\end{proof}

Thus the average excess of $w$ on the commutator subgroup measures exactly how much extra
weight went into the nonlinear coefficients. If all nonlinear coefficients are strictly
negative with arbitrary magnitudes, rescaling reduces to this normalization, and every
such nonnegative detector must count every nonidentity element of the commutator subgroup.

\begin{proposition}[When one Frobenius class suffices; \stU]\label{nab:prop:oneclass}
Let $D=[G,G]\ne1$. Call the set of generators of all conjugates of one cyclic subgroup a
\emph{rational class}. A test $f$ as in \cref{nab:thm:minimax}, supported on one
nonidentity rational class and with strictly negative coefficients on all nonlinear
irreducible characters, exists if and only if $D\setminus\set1$ is one rational class. This
holds if and only if $D$ is the additive group of a finite field $\F_{p^d}$ and, under a
suitable identification, the image $\Gamma$ of $G$ acting on $D$ by conjugation consists
of field multiplications with $\Gamma\,\F_p^\times=\F_{p^d}^\times$. The same statements
hold for ordinary conjugacy classes, with the condition $\Gamma=\F_{p^d}^\times$.
\end{proposition}
\begin{proof}
A class, rational or ordinary, containing an element of $D$ is contained in $D$, and
\cref{nab:thm:minimax}(b) forces it to contain $D\setminus\set1$; conversely use $f_*$.
For the structure, suppose $D\setminus\set1$ is one rational class. All nonidentity
elements of $D$ have the same order, necessarily a prime $p$, so $D$ is a $p$-group. Its
centre is nontrivial, characteristic, and stable under conjugation and under powers prime
to $p$, so transitivity forces $D=Z(D)$, and $D$ is an elementary abelian group $V$. The
conjugation image $\Gamma$ is abelian because $G/D$ is abelian. The group generated by
$\Gamma$ and the scalars $\F_p^\times$ acts transitively on $V\setminus\set0$, so $V$ is an
irreducible $\F_p[\Gamma]$-module. The commutative algebra generated by $\Gamma$ in
$\End_{\F_p}(V)$ is a finite
field, since every nonzero element has an invariant kernel and is therefore invertible,
with inverse in the algebra by finiteness; irreducibility makes $V$ one-dimensional over
it. The transitivity statements give the scalar conditions, which are also sufficient.
\end{proof}

For the group $U_3(\F_p)$ of upper unitriangular $3\times3$ matrices the commutator
subgroup is the centre, the nonzero centre is one rational class consisting of $p-1$
ordinary classes, and its density is $(p-1)/p^3$. For $p>2$ one ordinary class cannot meet
the sign and coverage requirements of \cref{nab:thm:minimax}.

\paragraph{Partial detection under a support budget.}
If the test may discard some new characters, the answer is a finite optimization over
normal subgroups. Define the mass of the active characters of $f$ by
\[
 \mathcal W(f)=\sum_{\chi\ \mathrm{new},\ a_\chi<0}\chi(1)^2 .
\]

\begin{theorem}[Support budget; \stU]\label{nab:thm:budget}
Among the functions $f$ of \cref{nab:thm:minimax}, for every $0\le B\le g$,
\[
 \max_{\abs{\supp f}\le B}\mathcal W(f)
 =\max_{\substack{K\trianglelefteq G,\ K\subseteq N\\ n-\abs K\le B}}
   \Bigl(\frac g{\abs K}-\frac gn\Bigr).
\]
For each feasible $K$ the indicator $f_K=\nabone_N-\nabone_K$ attains the value attached to
$K$ with exactly $n-\abs K$ elements in its support, and its new coefficients are
$-\abs K\chi(1)/g$ if $K\subseteq\ker\chi$ and $0$ otherwise.
\end{theorem}
\begin{proof}
By \cref{nab:thm:minimax}(a) the zero set $K$ of $f$ in $N$ is normal in $G$, $f$ factors
through $G/K$, its support on $N$ has $n-\abs K$ elements, and its active characters are
characters of $G/K$ not inflated from $G/N$. The sum of their squared degrees is at most
$g/\abs K-g/n$. Conversely, averaging a representation over the normal subgroup $K$
projects onto its $K$-invariants, which is the whole representation or zero; this gives the
coefficients of $f_K$ and the attainment.
\end{proof}

By \cref{nab:thm:nonsplit} applied to $E^K$, the arithmetic test attached to $f_K$ sees
exactly the Riemann hypothesis for $\zeta_{E^K}$. Two cases show how rigid the budget is.
For $G=S_4$ and $N=A_4$ the normal subgroups of $G$ inside $N$ are $A_4$, $V_4$ and $1$, so
the maximum of $\mathcal W$ is $0$ for $B<8$, $4$ for $8\le B<11$ and $22$ for $B\ge11$; at
the middle value only the two-dimensional representation inflated from $S_4/V_4\simeq S_3$
occurs, and the indicator of the eight $3$-cycles attains the bound
(Example~\ref{nab:ex:nonsplit}(b)). For a nonabelian simple group $G=N$, one active character
already forces support on all $g-1$ nonidentity elements; for $A_5$ the threshold $59$ is
all-or-nothing, even if the aim is a single irreducible character.

\paragraph{The cost of allowing both signs.}
A sparse observable might trade a strictly one-sided response for a small contribution of
the opposite sign. Let $G$ be nonabelian simple, and for $\chi\ne\nabone$ put
\[
 \ell_\chi(h)=1-\frac{\Re\chi(h)}{\chi(1)},\qquad
 P_\pm(f)=\sum_{\chi\ne\nabone,\ \pm a_\chi>0}\abs{a_\chi}\chi(1),\qquad
 \kappa(G)=\min_{h\ne1}\frac{\min_{\chi\ne\nabone}\ell_\chi(h)}{\max_{\chi\ne\nabone}\ell_\chi(h)} .
\]
Every nontrivial irreducible representation is faithful, so $\ell_\chi(h)>0$ for $h\ne1$.

\begin{proposition}[\stU]\label{nab:prop:kappa}
If a nonzero real inversion-symmetric class function $0\le f\le1$ on a nonabelian simple
group $G$ vanishes at $1$ and at some $h\ne1$, then $P_+(f)\ge\kappa(G)P_-(f)$. If
$\kappa(G)<1$ the bound is attained.
\end{proposition}
\begin{proof}
The equation $f(h)-f(1)=0$, written as in \eqref{nab:eq:identity} over all
$\chi\ne\nabone$, equates the negative and positive masses weighted by their losses
$\ell_\chi(h)$. For the attainment choose $h,\chi,\psi$ realizing the minimum ratio and
take a positive multiple of $f(u)=\ell_\chi(u)-\kappa(G)\ell_\psi(u)$. It is nonnegative by
the definition of $\kappa(G)$ and vanishes at $h$ and at $1$. The two character pairs,
with complex conjugates included as needed, are disjoint because their normalized real
traces differ at $h$; hence $P_+/P_-=\kappa(G)$. Scale $f$ to have maximum one.
\end{proof}

For $G=A_5$ the class-by-class ratios behind $\kappa(G)$ and the complete rational frontier of
$P_+/P_-$ under a support budget are computed in \cref{nab:prop:a5frontier}
(\cref{grp:sec:a5}).

\paragraph{Binary tests.}
At an unramified prime one often records a yes/no event, such as whether a polynomial
stays irreducible. Let $N\ne1$ be normal in $G$, $\pi:G\to Q=G/N$, $q=[G:N]$, $n=\abs N$.
For a real function $f$ on $G$ with $f(u^{-1})=f(u)$ let $A$ be the matrix
$A_{uv}=f(u^{-1}v)$ on $\R^G$, with standard basis $(e_u)_{u\in G}$, and let
$V\subset\R^G$ be the space of vectors whose coordinates sum to zero on each coset of $N$. Left translation by $N$ acts transitively on
each coset, so the coset blocks of $A$ have constant row and column sums and $A$ preserves
$V$. If $f$ is a class function, $A$ acts on the isotypic subspace of an irreducible
character $\chi$ by the scalar $\abs G\langle f,\chi\rangle/\chi(1)$, and the characters
occurring in $V_\C$ are exactly the new ones. Thus $A|_V<0$, respectively $A|_V\le0$, says
that every new coefficient is negative, respectively nonpositive.

\begin{proposition}[Integral rigidity; \stU]\label{nab:prop:rigidity}
Suppose $f$ takes integer values and $A|_V$ is negative definite. Then
\[
 \sum_{u\in N}\bigl(f(u)-f(1)\bigr)\ge n-1,
\]
with equality if and only if $f=k\circ\pi-\nabone_{\set1}$ for an integer-valued
inversion-symmetric function $k$ on $Q$. In particular, for $f$ with values in $\set{0,1}$,
$A|_V$ is negative definite if and only if
\[
 f=\nabone_{\pi^{-1}(D)}-\nabone_{\set1},\qquad 1\in D=D^{-1}\subseteq Q .
\]
\end{proposition}
\begin{proof}
$T=-A|_V$ preserves the full lattice $V\cap\Z^G$ and is positive definite, so $\det T$ is a
positive integer. On the coset-constant subspace the trace of $A$ is $q\sum_{u\in N}f(u)$,
using normalized coset indicators as an orthonormal basis, and the full trace is $qnf(1)$.
Hence
\[
 1\le(\det T)^{1/(q(n-1))}\le\frac{\tr T}{q(n-1)}=\frac{\sum_{u\in N}(f(u)-f(1))}{n-1}.
\]
Equality forces all eigenvalues of $T$ to equal one. Then $(A+I)(e_u-e_v)=0$ for $u,v$ in
the same coset, and reading the row of the identity gives the factorization through the
quotient. Conversely the convolution matrix of a function pulled back from $Q$ annihilates
$V$. In the binary case the sum is at most $n-1$, so equality is forced and $f(1)=0$.
\end{proof}

For a binary class function and $N=[G,G]$ this recovers \cref{nab:thm:minimax}: all
strictly negative nonlinear coefficients are forced to equal $-\chi(1)/\abs G$, and the
minimal support is $N\setminus\set1$, of density $(n-1)/\abs G$. If zero coefficients are
allowed the answer changes: a quotient, a normal subgroup of that quotient and a sign
character can all occur.

\begin{theorem}[Binary tests with zero coefficients; \stU]\label{nab:thm:binary}
Let $f:G\to\set{0,1}$ be a class function with $f(x^{-1})=f(x)$, and call $f$
\emph{admissible} if $A|_V\le0$. Functions constant on the cosets of $N$ are admissible and
annihilate $V$. Every other admissible $f$ factors through a quotient $\bar G=G/K$ for some
$K\triangleleft G$ strictly contained in $N$. Write $\bar N=N/K$ and $\bar\pi:\bar G\to Q$.
There exist a normal subgroup $J\triangleleft\bar G$ with $J\cap\bar N=1$, a
conjugation-invariant homomorphism $\varepsilon:J\to\set{\pm1}$, and a binary class function
$b$ on $Q$, invariant under inversion, such that
\begin{equation}\label{nab:eq:binary}
 f(x)=b(\bar\pi(x))-\varepsilon(x)\nabone_J(x)\qquad(x\in\bar G),
\end{equation}
where $\varepsilon\nabone_J$ is extended by zero and
$b(\bar\pi(j))=(1+\varepsilon(j))/2$ for $j\in J$. Conversely every such function is
admissible. On $V$ its only possible nonzero eigenvalue is $-\abs K\abs J$, with
multiplicity $(\abs{G/K}-\abs{G/N})/\abs J$. A new character has a nonzero coefficient
exactly when $K$ acts trivially and $J$ acts on the resulting representation of $\bar G$ by
the scalar character $\varepsilon$.
\end{theorem}
\begin{proof}
For $u,v$ in the same coset of $N$, the negative of the quadratic form of $A$ on $e_u-e_v$
is $2(f(u^{-1}v)-f(1))$. If it vanishes, then $A(e_u-e_v)=0$, by semidefiniteness and
invariance of $V$. If $f(1)=1$ all these differences vanish, and $f$ is constant on each
coset. Otherwise $f(1)=0$. The elements $h\in N$ with $f(h)=0$ are exactly those with
$Ae_h=Ae_1$. By translation invariance they form the intersection of $N$ with the
stabilizer of a column, hence a subgroup $K$, which is normal because $f$ is a class
function; equality of columns and translation show that $f$ is constant on the cosets of
$K$. If $K=N$ the first case applies. Otherwise pass to $\bar G=G/K$: the operator on the
$K$-coset-constant subspace is $\abs K$ times the operator of the quotient, and it vanishes
on the orthogonal complement. We now work in $\bar G$ and write $\nu=\abs{\bar N}\ge2$.

In the identity coset $f$ is $1$ off the identity and $0$ at the identity. Split $V$ into
the spaces $W_C$ of vectors of sum zero on the individual cosets $C$ of $\bar N$. For
$B=-A|_V$ every diagonal block is the identity of $W_C$, and each off-diagonal block is a
contraction, by positivity of the compression to two blocks with identity diagonal. For
cosets $C,D$ the binary matrix of $A$ between them is regular with row and column sum
$r=\sum_{x\in C^{-1}D}f(x)$. Its squared Frobenius norm is $\nu r$, and the constant
direction accounts for $r^2$. At most $\nu-1$ remaining singular values are nonzero and
each is at most one, so $r(\nu-r)\le\nu-1$, and $r\in\set{0,1,\nu-1,\nu}$.

A coset on which $f$ is constant gives a zero block. Any other coset of $\bar N$, called
active, has one exceptional element $s_\alpha$ over $\alpha\in Q$: a missing element among
ones, or a present element among zeros. Give it the sign $\varepsilon_\alpha=+1$ in the
first case and $-1$ in the second; when $\nu=2$ always choose the missing element. On the
sum-zero spaces the corresponding block of $B$ is $\varepsilon_\alpha$ times the
permutation identifying the two cosets by right multiplication by $s_\alpha$. The identity
coset is active, with $s_1=1$ and positive sign.

Write the positive block matrix as a Gram matrix $B_{CD}=U_C^*U_D$ with isometries $U_C$.
A block that is an isometry identifies the two image subspaces, since orthogonal
projection onto one image preserves the norms of all vectors of the other and the
dimensions agree. Consequently, if $\alpha,\alpha'\in Q$ are active, the blocks between
the cosets $1,\alpha,\alpha\alpha'$ satisfy
$B_{1,\alpha\alpha'}=B_{1,\alpha}B_{\alpha,\alpha\alpha'}$. This is an isometry, so
$\alpha\alpha'$ is active, and the active cosets form a subgroup of $Q$. In the last
identity the product of the exceptional matchings is right multiplication by
$s_\alpha s_{\alpha'}$. On the sum-zero space a permutation is determined by its
restriction, and for $\nu\ge3$ no permutation acts as minus the identity there, since it
would have to interchange each pair of distinct coordinates. Thus
$s_{\alpha\alpha'}=s_\alpha s_{\alpha'}$ and
$\varepsilon_{\alpha\alpha'}=\varepsilon_\alpha\varepsilon_{\alpha'}$. For $\nu=2$ all
chosen signs are positive and faithfulness gives the same conclusions. Conjugation
preserves the exceptional element and the sign of every active coset. The exceptional
elements therefore form a normal subgroup $J$ with $J\cap\bar N=1$, and their signs give
an invariant character. Filling in the constant background on each coset gives
\eqref{nab:eq:binary}.

Conversely, the first term of \eqref{nab:eq:binary} annihilates $V$. On the $K$-invariant
part of $V$ the operator $-A$ acts as $\abs K\sum_{j\in J}\varepsilon(j)R_j$, where
$(R_jv)(x)=v(xj)$; this is $\abs K\abs J$ times the orthogonal projection onto the
$\varepsilon$-isotypic subspace for $J$. On all $K$-invariant vectors this projection has
rank $\abs{G/K}/\abs J$, by the regular action of $J$ on each of its cosets, and on
$N$-invariant vectors the rank is $\abs{G/N}/\abs J$, since $J$ embeds in $Q$. Subtraction
gives the multiplicity. The projection is central because $\varepsilon$ is invariant and
$J$ is normal, which gives the criterion on irreducible characters.
\end{proof}

\begin{corollary}[\stU]\label{nab:cor:binarysn}
For $G=S_n$, $n\ge5$, and $N=A_n$, a binary test all of whose nonlinear coefficients are
nonpositive either annihilates all nonlinear characters or detects all of them, with the
same eigenvalue $-1$ on $V$.
\end{corollary}
\begin{proof}
$A_n$ is simple, so $K=1$ in the nonconstant case. Normal subgroups with trivial
intersection commute, hence $J\subseteq C_{S_n}(A_n)=1$: an element centralizing every
$3$-cycle preserves the support of every $3$-cycle, hence every point.
\end{proof}

\begin{example}[The smallest selective event]\label{nab:ex:twisted}
Let $L$ be the splitting field of $X^3-X-1$, with group $S_3$, and $E=L(i)$, so
$G=S_3\times C_2$. Let $\sigma$ be the sign character, $\rho$ the standard character of
degree two and $\eta$ the nontrivial character of $C_2$. For $p\nmid46$ let $f(p)=1$ if
either $p\equiv1\pmod4$ and $X^3-X-1$ is irreducible modulo $p$, or $p\equiv3\pmod4$ and it
has three distinct roots modulo $p$; let $f(p)=0$ otherwise. The corresponding class
function is
\[
 f=\frac{1+\sigma}4+\frac{(1+\sigma)\eta}{12}-\frac{\rho\eta}3 ,
\]
with mean $1/4$ and no $\rho$ term. It is the unique inversion-invariant binary class
function of least support having a negative coefficient on $\rho\eta$ and zero coefficient
on $\rho$; its support has three of the twelve elements. Indeed, here $N=A_3\times1$ has
order three; for a selective test in \cref{nab:thm:binary}, $K=1$ and $J\ne1$; since $J$
centralizes $N$ and meets it trivially, $J=1\times C_2$; to select $\rho\eta$ the sign must
be $\eta$; the two active cosets then have two selected elements and one selected element,
and choosing zero on the two remaining cosets is the unique minimum. The one-sided
criterion for $L(s,\rho\eta)$ attached to this event is \cref{hec:thm:twistedcubic}.
\end{example}

\paragraph{Rotation classes and integrality.}
A single measured class may have character coefficients of both signs and still detect
zeros, because zero multiplicities are integers. Let $G=A\rtimes\set{1,\iota}$ with $A$ a
nontrivial abelian group of odd order $n$ and $\iota a\iota^{-1}=a^{-1}$; the dihedral
group $D_{2n}$ is the case $A=C_n$. The nonlinear irreducible representations of $G$ are
the inductions of the nontrivial characters $\psi$ of $A$, indexed by the pairs
$\set{\psi,\psi^{-1}}$: they are irreducible by the index-two character calculation, and
their squared degrees together with the two linear characters sum to $2n$. Write a
nonidentity rotation class as $[a]=\set{a,a^{-1}}$. For a vector $m=(m_\psi)$ with
$m_{\nabone}=0$ and $m_\psi=m_{\psi^{-1}}$ put
\[
 T_m(a)=\sum_{\psi\in\widehat A}m_\psi\psi(a) .
\]

\begin{theorem}[Rotation classes; \stU]\label{nab:thm:rotation}
A collection $\mathcal S$ of nonidentity rotation classes detects every nonzero vector $m$
of nonnegative integers, in the sense that $T_m(a)\ne0$ for some $[a]\in\mathcal S$, if and
only if $\mathcal S$ contains the class of a generator of every nontrivial cyclic subgroup
of $A$. The least number of classes is the number of nontrivial cyclic subgroups of $A$;
for $D_{2n}$ it is $\tau(n)-1$, with $\tau$ the divisor function, and for prime $n$ every
single nonidentity rotation class suffices.
\end{theorem}
\begin{proof}
For $u$ prime to the order of $a$, the automorphism $\zeta\mapsto\zeta^u$ of the
cyclotomic field carries $T_m(a)$ to $T_m(a^u)$, since the $m_\psi$ are integers. If every
cyclic subgroup has a generator among the measured classes, vanishing there forces $T_m$
to vanish on $A\setminus\set1$; Fourier inversion on $A$ makes all $m_\psi$ equal, and
$m_{\nabone}=0$ gives $m=0$.

Conversely, let $C$ be a nontrivial cyclic subgroup with no measured generator, $R$ its set
of generators, and
\[
 m_\psi=\abs R-\sum_{a\in R}\overline{\psi(a)} .
\]
The inner sum is a sum of roots of unity invariant under every automorphism of the
cyclotomic field of order $\abs C$, hence a rational integer, and it is at most $\abs R$.
So the $m_\psi$ are symmetric nonnegative integers with $m_{\nabone}=0$, and by character
orthogonality
\[
 T_m(b)=n\abs R\,\nabone_{\set1}(b)-n\,\nabone_R(b) .
\]
Thus $m$ is nonzero and invisible at every measured class. Distinct cyclic subgroups have
disjoint sets of generators, which gives the minimum. For cyclic $A$ the inner sum is a
Ramanujan sum; in $C_9$ the class of a generator alone misses the multiplicities
$(3,3,0,3)$ at the pairs $\set{\pm j}$, $j=1,2,3,4$, and an additional class of order three
gives full detection.
\end{proof}

\begin{corollary}[\stG; \stR]\label{nab:cor:rotation}
Let $E/\Q$ be Galois with group $G=A\rtimes\set{1,\iota}$ as above, and let $L_\psi$ be the
$L$-functions of its nonlinear irreducible characters; each is the Hecke $L$-function of a
nontrivial character of finite order of the quadratic field $E^A$. For a rotation class
$[a]$ put
\[
 S_a(x)=\sum_{\substack{p\le x,\ p\notin S\\g_p\in[a]}}\log p\,\Bigl(1-\frac px\Bigr).
\]
Choose one generator class for each nontrivial cyclic subgroup of $A$. Given
\cref{src:thm:dgrh}, all $L_\psi$ satisfy the Riemann hypothesis if and only if for every
selected $[a]$ and every $\eps>0$,
\[
 S_a(x)=\frac x{2n}+O_{E,\eps}(x^{1/2+\eps}).
\]
Under the Riemann hypothesis the error is $O_E(\sqrt x)$. If the nonlinear factors are
known to have no real zeros in $(1/2,1)$, an upper bound for each centered sum suffices.
\end{corollary}
\begin{proof}
Inserting all prime powers into $S_a$ costs $O_E(\sqrt x)$. The indicator of $[a]$ has
coefficient $1/n$ on the trivial character and $(\psi(a)+\overline{\psi(a)})/n$ on the
nonlinear character attached to $\set{\psi,\psi^{-1}}$. The pole at $1$ contributes
$x/(2n)$, and at a point $\rho$ in the right half of the strip the Mellin transform of the
centered sum against $x^{-s-1}\,\dd x$ has residue $-T_m(a)/(n\rho(\rho+1))$, where $m_\psi$
is the order of $L_\psi$ at $\rho$; \cref{src:thm:dgrh} takes care of the two linear terms.
The two-sided estimate makes the transform holomorphic to the right of $1/2$, so
\cref{nab:thm:rotation} excludes every zero. For the converse apply the explicit formula
for the Riesz mean to each fixed factor (\cref{sec:hec}). If real zeros are already excluded, the transform
of a nonnegative majorant minus the centered sum has no real singularity to the right of
the proposed exponent, and \cref{ana:thm:onesided} gives the same holomorphy from the
one-sided estimate. No hypothesis of linear independence of zero ordinates, simplicity or
disjointness of zero sets is used.
\end{proof}


\subsection{What fixed congruence weights can see}\label{nab:sec:periodic}
\begingroup\raggedright
\sources{grh_method_transfers_2026-11-25; grh-arithmetic-transfers; grh_sylow_transfer_2026-11-25; grh_relative_prime_races_2026-11-25; grh_deficit_transfer_2026-11-25; grh_cyclic_generator_transfer_2026-11-25; psl2_grh_transfer}
\endgroup

The prime weights produced by the source proof are functions of a residue class of $p$
(\cref{src:sec:mechanism}). The following theorem computes exactly how well such a
function can approximate a class function of Frobenius. It is an application of character
theory and of the unconditional Chebotarev theorem for a fixed field.

Fix an integer $q\ge1$ and put
\[
 A_q=(\Z/q\Z)^\times,\qquad B_q=E\cap\Q(\zeta_q),\qquad H_q=\Gal(E/B_q).
\]
Since $B_q/\Q$ is abelian, $H_q$ is normal and contains $[G,G]$. Assume that $S$ also
contains the primes dividing $q$. The restriction of $b\in A_q=\Gal(\Q(\zeta_q)/\Q)$ to
$B_q$ determines a coset of $H_q$ in $G$. For a class function $f:G\to\C$ define
\begin{equation}\label{nab:eq:periodic-best}
 m_q(b)=\frac1{\abs{H_q}}\sum_{\substack{g\in G\\g|_{B_q}=b|_{B_q}}}f(g),\qquad
 \norm f_G^2=\frac1{\abs G}\sum_{g\in G}\abs{f(g)}^2,\qquad
 \norm a_{A_q}^2=\frac1{\varphi(q)}\sum_{b\in A_q}\abs{a(b)}^2 .
\end{equation}

\begin{theorem}[Optimal approximation by a fixed modulus; \stU]\label{nab:thm:periodic}
For every class function $f:G\to\C$ and every $a:A_q\to\C$,
\[
 \lim_{X\to\infty}\frac1{\pi(X)}\sum_{\substack{p\le X\\p\notin S}}\abs{f(g_p)-a(p\bmod q)}^2
 =\lim_{X\to\infty}\frac1X\sum_{\substack{p\le X\\p\notin S}}(\log p)\abs{f(g_p)-a(p\bmod q)}^2
 =\mathcal E_q(f,a),
\]
where
\begin{equation}\label{nab:eq:mse}
 \mathcal E_q(f,a)=\norm f_G^2-\norm{m_q}_{A_q}^2+\norm{a-m_q}_{A_q}^2 .
\end{equation}
Consequently $m_q$ is the unique optimal approximation on $A_q$, and if
$f=\sum_{\chi\in\nabIrr(G)}c_\chi\chi$, the minimum error is
\begin{equation}\label{nab:eq:energy}
 \min_a\mathcal E_q(f,a)=\norm f_G^2-\norm{m_q}_{A_q}^2
 =\sum_{\substack{\chi\in\nabIrr(G)\\H_q\not\subseteq\ker\chi}}\abs{c_\chi}^2 .
\end{equation}
\end{theorem}
\begin{proof}
Put $M=E\,\Q(\zeta_q)$. Restriction identifies $\Gal(M/\Q)$ with the fibre product
$\Gamma=\set{(g,b)\in G\times A_q:g|_{B_q}=b|_{B_q}}$, whose projection to $A_q$ is
surjective with fibres of size $\abs{H_q}$. The function
$F(g,b)=\abs{f(g)-a(b)}^2$ is constant on the conjugacy classes of $\Gamma$, because $f$ is
a class function and $A_q$ is abelian. Chebotarev's density theorem, summed over the
finitely many classes, gives
\[
 \lim_{X\to\infty}\frac1{\pi(X)}\sum_{\substack{p\le X\\p\notin S}}F(g_p,p\bmod q)
 =\frac1{\varphi(q)\abs{H_q}}\sum_{b\in A_q}\sum_{g|_{B_q}=b|_{B_q}}\abs{f(g)-a(b)}^2 .
\]
For fixed $b$ the mean of $f$ over the indicated coset is $m_q(b)$, and expanding the
square gives
\[
 \frac1{\abs{H_q}}\sum_{g|_{B_q}=b|_{B_q}}\abs{f(g)-a(b)}^2
 =\frac1{\abs{H_q}}\sum_{g|_{B_q}=b|_{B_q}}\abs{f(g)-m_q(b)}^2+\abs{m_q(b)-a(b)}^2 .
\]
Average over $b$; every $g\in G$ is counted equally often because
$A_q\to\Gal(B_q/\Q)$ is surjective. This proves \eqref{nab:eq:mse} and the uniqueness of
the minimum. The logarithmic weight does not change the limit: if $U(X)$ is the
unweighted sum, then $U(X)=\mathcal E_q(f,a)\Li(X)+o(X/\log X)$ and partial summation gives
$U(X)\log X-\int_2^XU(t)t^{-1}\,\dd t=\mathcal E_q(f,a)X+o(X)$.

For the character expression, let $\varrho$ be an irreducible representation of $G$ with
character $\chi$, and $P=\abs{H_q}^{-1}\sum_{h\in H_q}\varrho(h)$ the projection onto the
$H_q$-invariants. Because $H_q$ is normal this subspace is $G$-invariant, hence zero or
everything, and
\[
 \frac1{\abs{H_q}}\sum_{h\in H_q}\chi(gh)=\tr(\varrho(g)P)=
 \begin{cases}\chi(g),&H_q\subseteq\ker\chi,\\0,&H_q\not\subseteq\ker\chi.\end{cases}
\]
Orthogonality of irreducible characters shows that the squared norm lost by the coset
average is the sum in \eqref{nab:eq:energy}.
\end{proof}

\begin{corollary}[Nonlinear traces are invisible to fixed congruences; \stU]\label{nab:cor:nonlinear}
If $\chi$ is an irreducible character of $G$ of degree greater than one, then for every
fixed $q$ and every $a:A_q\to\C$,
\[
 \mathcal E_q(\chi,a)=1+\frac1{\varphi(q)}\sum_{b\in A_q}\abs{a(b)}^2 .
\]
The unique optimal periodic approximation is zero, and the minimum mean-square error is
exactly one. For the trace divided by the degree $d$ it is $1/d^2$.
\end{corollary}
\begin{proof}
$G/H_q$ is abelian, so a nonlinear irreducible representation is not trivial on $H_q$;
hence $m_q=0$, and $\norm\chi_G^2=1$.
\end{proof}

Let $E^{\mathrm{ab}}=E^{[G,G]}$. By the Kronecker--Weber theorem it lies in
$\Q(\zeta_Q)$ for some $Q$, and for every multiple $q$ of $Q$ one has
$B_q=E^{\mathrm{ab}}$ and $H_q=[G,G]$. For such $q$,
\begin{equation}\label{nab:eq:allmoduli}
 \min_a\mathcal E_q(f,a)=\sum_{\chi\in\nabIrr(G),\ \chi(1)>1}\abs{c_\chi}^2 ,
\end{equation}
and this is also the infimum over all fixed moduli and all periodic approximations; for
smaller moduli one may in addition lose linear characters whose abelian fields are not in
the chosen cyclotomic field. A class function is exactly recoverable in prime mean square
from some fixed residue modulus if and only if it factors through the abelianization of
$G$. Every function of a fixed residue class is a finite linear combination of Dirichlet
characters. Therefore none of the weights of this section that involve a nonlinear
character can be written, even in mean square, as a fixed finite combination of
Dirichlet-character weights: this applies to $W^{\mathrm{Syl}}$ and $W^{\mathrm{cyc}}$ of
\cref{nab:thm:sylow}, whose values at the identity and at a nonidentity element of a Sylow
subgroup of $[G,G]$ differ although both elements have the same image in every abelian
quotient, to the compressed weights of \cref{nab:thm:tower}, to
$\nabone_{[G,G]\setminus\set1}$, and to the detectors of \cref{sec:grp}, whose minimal
errors are collected in the table after \cref{nab:cor:doublecoset}. An extension of the source proof to these weights
must therefore use arithmetic information beyond fixed congruences; substituting a
character-table entry for a Dirichlet symbol in the final inequality is not enough.

\begin{remark}\label{nab:rem:fixed}
The quantifier \emph{fixed} is essential. The theorem fixes $q$ and $a$ and then lets
$X\to\infty$. It makes no assertion when $q=q(X)$ grows or when the residue weight is
chosen from the primes up to $X$: if $q>X$ is prime, each prime $p\le X$ occupies its own
residue class and any finitely many traces can be interpolated exactly. Such an
interpolation supplies neither analytic control nor an arithmetic construction. Nor does
the theorem prove that Artin or Hecke GRH is logically independent of Dirichlet GRH.
\end{remark}

For permutation characters the minimum has a closed form, and the optimal predictor is
the weight subtracted in \cref{nab:thm:missing}.

\begin{corollary}[Euler coefficients of a number field; \stU]\label{nab:cor:doublecoset}
Let $K$ be a number field with normal closure $E$, $U=\Gal(E/K)$ and $F=E^{U[G,G]}$. Over
all complex-valued periodic functions $b$ of $p$,
\[
 \min_b\lim_{x\to\infty}\frac1{\pi(x)}\sum_{p\le x}\abs{a_K(p)-b(p)}^2
 =\abs{U\backslash G/U}-[F:\Q],
\]
and an optimal $b$ equals $a_F(p)$ outside a finite set.
\end{corollary}
\begin{proof}
Apply \cref{nab:thm:periodic} to $u=\Ind_U^G\nabone$ with a modulus $q$ for which
$H_q=[G,G]$. The coset average of $u$ is the character of the $[G,G]$-invariants of the
permutation representation on $G/U$; a basis is indexed by the $[G,G]$-orbits, which form
the $G$-set $G/U[G,G]$, whose character is $a_F$ on Frobenius elements. This character is
$[G:U[G,G]]$ on the normal subgroup $U[G,G]$ and zero elsewhere, so its squared norm is
$[G:U[G,G]]=[F:\Q]$. The squared norm of $u$ is the number of $G$-orbits on $(G/U)^2$ by
Burnside's lemma, that is the number of double cosets.
\end{proof}

For $K=\Q(\sqrt[3]2)$ the best approximation is the constant $1$: the values $3,1,0$ of
$a_K$ occur with densities $1/6,1/2,1/3$ and the squared error is $4/6+1/3=1$; knowing
the sign of Frobenius does not improve it
(\artifact{grh-arithmetic-transfers/checks/rigidity.py}).

The explicit constants obtained in Part~II by reading off \eqref{nab:eq:energy},
\cref{nab:cor:nonlinear} or \cref{nab:prop:error} for particular groups are collected
below; the statements stay with the detectors they concern. The second column is the
weight, the third the least limiting error over all functions of $p$ modulo a fixed
modulus (mean square, except in the row marked ``binary'', where it is the density of
wrong decisions).
\begin{center}\small
\begin{tabular}{@{}llll@{}}\toprule
group or field & weight & least error & statement\\\midrule
$A_5$ & $W_{21}$; $W_{16}$ & $7$; $4$ & \cref{grp:cor:a5-variance}\\
$\PSL_2(\F_q)$, $q$ even & root-count weight $W$ & $q$ & after \cref{grp:cor:psl-root}\\
$\F_5^2\rtimes\F_5^\times$, six quintics & $W$ of \cref{grp:thm:six} & $6$ &
 \cref{grp:cor:six-variance}\\
Kummer block $S$ & $\sum_\ell b_\ell V_\ell$ & $\sum_\ell b_\ell^2$ & \cref{grp:cor:periodic}\\
$\Q(\sqrt[3]2)$ & $a_K$ & $1$ & above\\
cubic of disc.\ $-23$ & trace $\chi_\rho$ & $1$ & \cref{hec:prop:meansq}(i)\\
cubic of disc.\ $-23$ & irreducible primes & $2/9$, or $1/9$ if $23\mid m$ &
 \cref{hec:prop:meansq}(ii)\\
cubic of disc.\ $-23$ & irreducible primes, binary & $1/6$ & \cref{hec:prop:meansq}(ii)\\
$D_{10}$, conductor $47$ & one rotation class & $3/25$ & \eqref{hec:eq:325}\\\bottomrule
\end{tabular}
\end{center}

The mean-square measure can be replaced by the frequency of wrong decisions.

\begin{proposition}[Best congruence test for a Frobenius event; \stU]\label{nab:prop:error}
Let $C\subseteq G$ be a union of conjugacy classes. Over all fixed moduli $m$ and all sets
of residue classes modulo $m$ used to predict whether $g_p\in C$, the minimum natural
density, among all primes, of the set of wrong decisions is
\[
 \frac1{\abs G}\sum_{T\in G/[G,G]}\min\set{\abs{T\cap C},\ \abs{T\setminus C}},
\]
and the minimum is attained.
\end{proposition}
\begin{proof}
Put $M=\Q(\zeta_m)$ and $B=E\cap M$. As in the proof of \cref{nab:thm:periodic}, for a
given residue class the Frobenius $g_p$ is equidistributed, for conjugation-invariant
tests, in a coset of $\Gal(E/B)$, which is a union of cosets of $[G,G]$. Even if the
individual $[G,G]$-coset is known, the majority decision on it makes at least the stated
number of errors. For $m$ with $E^{[G,G]}\subseteq M$ the residue class reveals exactly
this coset, and the majority value attains the bound.
\end{proof}

\begin{corollary}[\stU]\label{nab:cor:dirichletproduct}
For the event ``$p$ splits completely in a number field $K$'', that is $C=\set1$ in the
Galois group of the normal closure, the minimum error density is $0$ if $G$ is abelian and
$1/\abs G$ otherwise, attained in the second case by the rule that no prime splits
completely. Consequently, if $\zeta_K$ agrees in $\Re s>1$ with a finite product of integer
powers of Dirichlet $L$-functions up to finitely many Euler factors, then the normal
closure of $K$ is abelian over $\Q$.
\end{corollary}
\begin{proof}
The first statement is Proposition~\ref{nab:prop:error}. If $\zeta_K$ is such a product, then by
uniqueness of Dirichlet coefficients $a_K(p)$ is, for almost all $p$, a fixed linear
combination of values $\chi(p)$, hence periodic, and the condition $a_K(p)=[K:\Q]$ predicts
complete splitting with error density zero.
\end{proof}

Thus an argument for the Riemann hypothesis of the zeta function of a nonabelian normal
closure, or of a nonnormal field generating it, needs more than a finite product or
quotient of Dirichlet $L$-functions. The obstruction concerns that strategy only; zeros of
distinct functions can still be related by other means, as in
\cref{nab:sec:heilbronn}.

\subsection{Open questions}\label{nab:sec:open}

\begin{question}\label{nab:q:arithmetic}
Prove any one of the one-sided inequalities of this section for one nonabelian field.
Instances with very few Frobenius classes are: \textup{(ii)} of \cref{nab:thm:reftest} for
the two-class weight of Example~\ref{nab:ex:quintic}; the lower bound of
\cref{nab:thm:nonsplit} for the irreducible primes of $X^3-X-1$, for the quartic
nonresidue event of Example~\ref{nab:ex:nonsplit}\textup{(a)}, or for the two events of
Example~\ref{nab:ex:nonsplit}\textup{(b)}; the lower bound for $U_K$ in
\cref{nab:thm:missing} for a primitive sextic field with group $A_5$; and the bound on
$B_f$ in \cref{nab:cor:rootcount}. By \cref{nab:sec:periodic} no fixed combination
of congruence weights reproduces these weights. What is missing is an arithmetic object,
in the role of the lattice of integral forms of \cref{src:sec:mechanism}, whose local
divisibility at $p$ is governed by a nonabelian Frobenius weight, such as
$\nabone_{[G,G]\setminus\set1}(g_p)$ or the fixed-point defect
$\abs G\nabone_{\set1}(g_p)-\#\nabFix(g_p,G/C)=W_C(g_p)$, with sufficient positivity and
control of denominators, together with matching upper and lower covolume bounds that
cancel the same main term.
\end{question}

\begin{question}\label{nab:q:un}
\begin{enumerate}[label=\textup{(\alph*)}]
\item For $n\ge13$, does the character $U_n$ of Remark~\ref{nab:rem:un} lie in the positive rational
cone spanned by the characters of $S_n$ induced from linear characters of subgroups? A
family of identities \eqref{nab:eq:R} uniform in $n$ would give
$\ord_\rho\zeta_E\le(n-1)\ord_\rho\zeta_K$ in all degrees, where at present only
$2(n-H_n)$ is known.
\item\label{grp:q:vn} Is $V_n=\Ind_{C_2}^{S_n}\sigma-\grpstd_n$ a nonnegative rational combination
of characters induced from linear characters of subgroups, for every $n\ge4$?
By \cref{grp:rem:s5-induction} this would give the corresponding statement for
$U_n$ in every degree; \cref{grp:thm:s5-identity} settles $n=5$.
\end{enumerate}
\end{question}

\begin{question}\label{nab:q:lowdegree}
Can the detecting subfield $E^{\langle t\rangle}$ of degree $n!/2$ in Corollary~\ref{nab:cor:sn} be
replaced by subfields of lower degree without assuming Artin holomorphy?
\end{question}

\begin{question}\label{nab:q:nonreflection}
Which Galois groups that are not reflection groups admit decompositions into two
involutions and a linear character as in \cref{nab:thm:reflection}? Can local positivity
on larger monomial subgroups give a detector of comparable size for the alternating
groups, where the device of a linear character equal to $-1$ on the reflections is not
available?
\end{question}

\begin{question}\label{nab:q:line}
Lemma~\ref{nab:lem:heilbronn}(c) excludes zeros and poles of Artin $L$-functions off the
critical line once $\mathrm{RH}(\zeta_E)$ is known. Can the exact local orders of several
positive monomial detectors constrain the Heilbronn character at a zero on the critical
line enough to exclude negative coefficients there?
\end{question}
